\BourbakiExercisesSection

\begin{enumerate}
\def\labelenumi{\arabic{enumi}.}
\item
  Show that if E is a vector space \(\neq\{0\}\) over a field K with an infinity of elements, the set of generating systems of E is not inductive for the order relation \(\supset\) (form a decreasing sequence \((S_{n})\) of generating systems of E such that the intersection of the \(S_{n}\) is empty).
\item
  Let K be a field, L a subfield of K, I a set and \((x_{i})_{1\leqslant i\leqslant m}\) a free family of vectors in \(K_{s}^{I}\) such that the coordinates of each of the \(x_{i}\) belong to L. Let V be the vector subspace of \(K_{s}^{I}\) generated by the \(x_{i}\) ; show that there exists a finite subset J of I, with m elements, such that for every vector \(z=(\zeta_{\alpha})_{\alpha\in\mathbb{I}}\) of V, all the \(\zeta_{\alpha}\) belong to the subfield of K generated by L and the m elements \(\zeta_{\beta}\) where \(\beta\in J\) . (By applying Corollary 3 of no. 5, to the space \(L_{s}^{I}\) , show that it can be reduced to the case where there exist m indices \(\beta_{i}\in I\) ( \(1\leqslant i\leqslant m\) ) such that \(\Pr_{\beta_{i}}(x_{j})=\delta_{ij}\) ).
\item
  \begin{enumerate}
  \def\labelenumii{(\alph{enumii})}
  \tightlist
  \item
    Let G be a group and M an infinite subset of G. Show that if \(G'\) is the subgroup of G generated by M, then \(\text{Card}(G') = \text{Card}(M)\) .
  \end{enumerate}
\end{enumerate}

\begin{enumerate}
\def\labelenumi{(\alph{enumi})}
\setcounter{enumi}{1}
\item
  Let \(\mathbf{K}\) be a field and \(\mathbf{M}\) an infinite subset of \(\mathbf{K}\) . Show that if \(\mathbf{K}'\) is the subfield of \(\mathbf{K}\) generated by \(\mathbf{M}\) , then \(\operatorname{Card}(\mathbf{K}') = \operatorname{Card}(\mathbf{M})\) . (Consider two sequences \((\mathbf{A}_n)_{n \geqslant 0}, (\mathbf{P}_n)_{n \geqslant 0}\) of subsets of \(\mathbf{K}\) such that \(\mathbf{A}_0 = \mathbf{M}\) , \(\mathbf{P}_n\) is the multiplicative subgroup of \(\mathbf{K}^*\) generated by \(\mathbf{A}_n \cap \mathbf{K}^*\) and \(\mathbf{A}_{n+1}\) the additive subgroup of \(\mathbf{K}\) generated by \(\mathbf{P}_n\) ; then apply (a).)
\item
  Let K be a field and I an infinite set; show that \(\operatorname{Card}(\mathbf{K}) \leqslant \dim(\mathbf{K}_{s}^{\mathbf{I}})\) . (Reduce it to the case I = N and argue by reductio ad absurdum. Let B be a basis of \(K_{s}^{N}\) such that \(\operatorname{Card}(\mathbf{B}) < \operatorname{Card}(\mathbf{K})\) and let L be the subfield of K generated by the coordinates of all the elements of B; then \(\operatorname{Card}(\mathbf{L}) < \operatorname{Card}(\mathbf{K})\) . Form a sequence \((\xi_{n})_{n \geqslant 1}\) of elements of K such that for all n, \(\xi_{n}\) does not belong to the subfield of K generated by L and the \(\xi_{i}\) of index i \textless{} n; then apply Exercise 2 to obtain a contradiction, by considering the point \(x = (\xi_{n})\) of \(K_{s}^{N}\) .
\item
  Deduce from (c) that for every field K and every infinite set I,
\end{enumerate}

\[
\dim (\mathrm{K} _ {s} ^ {\mathrm{I}}) = (\operatorname{Card} (\mathrm{K})) ^ {\operatorname{Card} (\mathrm{I})}
\]

(``Erdös-Kaplansky theorem'').

4 Give an example of an infinite sequence \((\mathbf{F}_n)\) of vector subspaces of a vector space \(\mathbf{E}\) such that \(\operatorname{codim}\left(\bigcap_{n} \mathbf{F}_n\right) > \sum_{n} \operatorname{codim}(\mathbf{F}_n)\) . (Take \(\operatorname{codim}(\mathbf{F}_n) = 1\) for all \(n\) and use Exercise 3 \((d)\) .)

\begin{enumerate}
\def\labelenumi{\arabic{enumi}.}
\setcounter{enumi}{4}
\tightlist
\item
  Let \((\mathbf{H}_{\lambda})_{\lambda \in L}\) be a family of hyperplanes (passing through 0) of a vector space \(\mathbf{E}\) over a field \(\mathbf{K}\) , which is a covering of \(\mathbf{E}\) .
\end{enumerate}

\begin{enumerate}
\def\labelenumi{(\alph{enumi})}
\item
  Show that, if K is finite, then \(\operatorname{Card}(\mathbf{L}) \geqslant 1 + \operatorname{Card}(\mathbf{K})\) and, if K is infinite, \(\operatorname{Card}(\mathbf{L}) \geqslant \aleph_{0} = \operatorname{Card}(\mathbf{N})\) . (Prove by induction on r that if K has at least r elements, then E cannot be the union of r hyperplanes.) Show by examples that these inequalities cannot be improved without an additional hypothesis (when K is infinite, consider the space \(\mathbf{E} = \mathbf{K}_{s}^{(\mathbf{N})}\) ).
\item
  If E is finite-dimensional, show that \(\operatorname{Card}(\mathbf{L}) \geqslant 1 + \operatorname{Card}(\mathbf{K})\) . (Argue by induction on \(\dim(\mathbf{E})\) .) which is no longer true in infinite dim, consider \(K^{(N)}\)
\end{enumerate}

¶ 6. Let S be a non-empty finite subset of a vector K-space E, not containing 0. Suppose that there exists an integer \(k \geqslant 1\) with the following property: (*) for all \(x \in S\) , there is a partition of \(S = \{x\}\) into \(h\) free subsets with \(h \leqslant k\) .

For such a partition \((\mathbf{N}_i)_{1\leqslant i\leqslant h}\) and every sequence \((r_j)_{1\leqslant j\leqslant m}\) of integers of \(\{1,h\}\) a subspace \(\mathrm{F}_{r_1r_2\dots r_m}\) of \(\mathbf{E}\) is defined by induction on \(m\) as follows: \(\mathrm{F}_{r_1}\) is the subspace generated by \(\mathrm{N}_{r_1},\mathrm{F}_{r_1\dots r_p}\) the subspace generated by the intersection of \(\mathrm{F}_{r_1\dots r_{p-1}}\) and \(\mathrm{N}_{r_p}\) . Suppose that:

(**) for all \(x \in S\) and every partition \((N_i)_{1 \leqslant i \leqslant h}\) of \(S = \{x\}\) into at most \(k\) free subsets, there exists at least one sequence \((r_j)_{1 \leqslant j \leqslant m}\) such that

\[
x \notin \mathrm{F} _ {r _ {1} r _ {2} \dots r _ {m}}.
\]

\begin{enumerate}
\def\labelenumi{(\alph{enumi})}
\item
  Let \(n\) be the least of the integers \(m \geqslant 1\) for all possible choices of \(x \in S\) and let \((\mathbf{N}_i)_{1 \leqslant i \leqslant h}\) and \((r_j)_{1 \leqslant j \leqslant m}\) satisfy conditions (*) and (**). Show that of necessity \(n = 1\) . (Argue by reductio ad absurdum by assuming \(x \notin F_{r_1 r_2 \ldots r_n}\) and \(n > 1\) . Then necessarily \(x \in F_{r_n}\) . For simplicity write \(E_p = F_{r_1 r_2 \ldots r_p}\) for \(1 \leqslant p \leqslant n\) and \(x = \sum_i \alpha_i y_i\) with \(\alpha_i \in K\) and \(y_i \in N_{r_n}\) ; there is at least one \(y_i\) not belonging to \(E_{n-1}\) , denote it by \(y\) ; then write \(N = N_{r_n} - \{y\}\) , \(N_i' = N_i\) for \(1 \leqslant i \leqslant h\) and \(i \neq r_n\) , \(N_{r_n}' = N \cup \{x\}\) , so that \((\mathbf{N}_i')_{1 \leqslant i \leqslant h}\) is a partition of \(S - \{y\}\) consisting of free subsets. Define \(E_0' = E\) and for \(1 \leqslant p \leqslant n\) , \(E_p'\) as the subspace generated by the intersection of \(E_{p-1}'\) and \(N_{r_p}'\) . Now \(y \notin E_{n-1}\) ; let \(q\) be the least integer such that \(y \notin E_q\) ; show that \(E_q' \not\subset E_q\) . Let \(s\) be the least integer such that \(E_s' \not\subset E_s\) ; show that \(r_s = r_n\) ; deduce finally a contradiction from the relations \(y \in E_s, E_{s-1}' \subset E_{s-1}\) and \(E_s' \not\subset E_s\) .
\item
  Conclude from (a) that with the given hypotheses there exists a partition of S into h free subsets for an integer \(h \leqslant k\) .
\end{enumerate}

\begin{enumerate}
\def\labelenumi{\arabic{enumi}.}
\setcounter{enumi}{6}
\tightlist
\item
  Let S be a non-empty finite subset of a vector K-space E, not containing 0. Suppose that there exists an integer \(k \geqslant 1\) such that for every subset
\end{enumerate}

T of S, \(\text{Card}(T) \leqslant k.\text{rg}(T)\) . Show that there exists a partition of S into h free subsets for some integer \(h \leqslant k\) . (Prove that S satisfies condition ( \(^{**}\) ) of Exercise 6, arguing by induction on \(\text{Card}(S)\) and taking amongst all the subspaces \(F_{r_{1}\ldots r_{m}}\) one of those with the smallest possible dimension; if G is such a subspace and j an index for which \(\text{Card}(G \cap N_{j})\) is the least possible, prove that

\[
\begin{array}{r l} k. \operatorname{Card} (\mathrm{G} \cap \mathrm{N} _ {j}) & \leqslant \operatorname{Card} (\mathrm{G} \cap (\mathrm{S} - \{x \})) \\ & \leqslant \operatorname{Card} (\mathrm{G} \cap \mathrm{S}) \leqslant k. \operatorname{Card} (\mathrm{G} \cap \mathrm{N} _ {j}).) \end{array}
\]

\begin{enumerate}
\def\labelenumi{\arabic{enumi}.}
\setcounter{enumi}{7}
\item
  Let E be a vector space. Show that every endomorphism of E which is not a right divisor of 0 in End(E) is surjective (§ 2, Exercise 7).
\item
  Let K be a field.
\end{enumerate}

\begin{enumerate}
\def\labelenumi{(\alph{enumi})}
\item
  In the vector space \(\mathbf{E} = \mathrm{K}_s^{(\mathbf{Z})}\) , let \((e_n)_{n\in \mathbf{Z}}\) denote the canonical basis; let \(v\) be the automorphism of \(\mathbf{E}\) defined by the relations \(v(e_n) = e_{n + 1}\) for all \(n\in \mathbf{Z}\) . If \(u = 1_{\mathbb{E}} - v\) , show that \(u\) is an injective endomorphism of \(\mathbf{E}\) , but that \(\dim (\operatorname {Coker}(u)) = 1\) .
\item
  In the vector space \(\mathbf{E} = \mathrm{K}_s^{(\mathbb{N})}\) , let \((e_n)_{n\in \mathbb{N}}\) denote the canonical basis; let \(v\) be the endomorphism of \(\mathbf{E}\) defined by \(v(e_0) = e_0\) , \(v(e_n) = e_{n - 1} + e_n\) for \(n\geqslant 1\) . Show that \(v\) is an automorphism of \(\mathbf{E}\) and that, if \(u = 1_{\mathbb{E}} - v\) , \(u\) is a surjective endomorphism of \(\mathbf{E}\) such that \(\dim (\mathrm{Ker}(u)) = 1\) .
\item
  Let I be a non-empty set. Show that every automorphism u of the vector space \(E = K_{s}^{(I)}\) can be written as a difference u = v - w of two automorphisms of E, except when K has only two elements and I consists of a single element. (Note that it suffices to prove that one particular automorphism u is of the desired form; when K is a field with two elements, consider separately the cases where \(\text{Card}(I) = 2\) and \(\text{Card}(I) = 3\) .)
\end{enumerate}

¶ 10. Let E be a vector space over a field K. Show that every endomorphism u of E is an automorphism or can be written as u = v - w, where v and w are automorphisms of E. (Observe that u can always be replaced by \(s_{1} \circ u \circ s_{2}\) , where \(s_{1}\) and \(s_{2}\) are two automorphisms of E. Distinguish two cases, according to whether \(\dim(\text{Ker}(u)) \leqslant \dim(\text{Coker}(u))\) or \(\dim(\text{Ker}(u)) > \dim(\text{Coker}(u))\) ; when \(\text{rg}(u)\) is finite, then the first case always holds. When

\[
\dim (\operatorname{Ker} (u)) \leqslant \dim (\operatorname{Coker} (u)),
\]

it can be reduced to the case where \(\operatorname{Ker}(u) \cap \operatorname{Im}(u) = \{0\}\) ; if

\[
\dim (\operatorname{Ker} (u)) = \dim (\operatorname{Coker} (u)),
\]

it can also be assumed that \(\operatorname{Im}(u) + \operatorname{Ker}(u) = \mathbf{E}\) and then apply Exercise 9 (c). If \(\dim (\operatorname {Ker}(u)) < \dim (\operatorname {Coker}(u))\) , there is a supplementary subspace in \(\mathbf{E}\) of \(\operatorname {Ker}(u)\) of the form \(\mathbf{W} \oplus \operatorname{Im}(u)\) with \(\dim (\mathbf{W}) = \dim (\operatorname {Im}(u)) = \dim (\mathbf{E})\) ; apply the remark at the beginning with \(s_1(u(\operatorname {Im}(u))) = \operatorname{Im}(u)\) and use Exercise 9 (a) and (c). When \(\dim (\operatorname {Ker}(u)) > \dim (\operatorname {Coker}(u))\) , it can be reduced to the case where \(\operatorname{Ker}(u) \subset \operatorname{Im}(u)\) ; this time take \(s_1(\operatorname{Ker}(u)) = u(\operatorname{Ker}(u))\) and use the results of the preceding cases and Exercise 9.)

\begin{enumerate}
\def\labelenumi{\arabic{enumi}.}
\setcounter{enumi}{10}
\tightlist
\item
  Let E, F be two vector spaces over a field K and \(u: E \rightarrow F\) a linear mapping. Show that for every vector subspace V of F,
\end{enumerate}

\[
\dim (u ^ {- 1} (\mathbf {V})) = \dim (\mathbf {V} \cap \operatorname{Im} (u)) + \dim (\operatorname{Ker} (u)).
\]

\begin{enumerate}
\def\labelenumi{\arabic{enumi}.}
\setcounter{enumi}{11}
\tightlist
\item
  Let E, F be two vector spaces and u, v two linear mappings of E into F.
\end{enumerate}

\begin{enumerate}
\def\labelenumi{(\alph{enumi})}
\tightlist
\item
  Show that the following inequalities hold:
\end{enumerate}

\[
\begin{array}{c} \operatorname{rg} (v) \leqslant \operatorname{rg} (u + v) + \operatorname{rg} (u) \\ \operatorname{rg} (u + v) \leqslant \inf (\dim (\mathrm{E}), \dim (\mathrm{F}), \operatorname{rg} (u) + \operatorname{rg} (v)). \end{array}
\]

If E and F are finite-dimensional, show that \(\operatorname{rg}(u + v)\) can take any integer value satisfying the inequalities

\[
\left| \operatorname{rg} (u) - \operatorname{rg} (v) \right| \leqslant \operatorname{rg} (u + v) \leqslant \inf (\dim (\mathrm{E}), \dim (\mathrm{F}), \operatorname{rg} (u) + \operatorname{rg} (v)).
\]

\begin{enumerate}
\def\labelenumi{(\alph{enumi})}
\setcounter{enumi}{1}
\tightlist
\item
  Show that
\end{enumerate}

\[
\dim (\operatorname{Ker} (u + v)) \leqslant \dim (\operatorname{Ker} (u) \cap \operatorname{Ker} (v)) + \dim (\operatorname{Im} (u) \cap \operatorname{Im} (v)).
\]

\begin{enumerate}
\def\labelenumi{(\alph{enumi})}
\setcounter{enumi}{2}
\tightlist
\item
  Show that
\end{enumerate}

\[
\dim (\operatorname{Coker} (u + v)) \leqslant \dim (\operatorname{Coker} (u)) + \operatorname{rg} (v).
\]

(If \(\mathbf{V}\) is a supplementary subspace of \(\operatorname{Ker}(v)\) in \(\mathbf{E}\) , note that

\[
u (\operatorname{Ker} (v)) \subset \operatorname{Im} (u + v) \quad \text { and } \quad \dim (u (\mathbf {V})) \leqslant \operatorname{rg} (v).)
\]

\begin{enumerate}
\def\labelenumi{\arabic{enumi}.}
\setcounter{enumi}{12}
\tightlist
\item
  Let E, F, G be three vector spaces and \(u: \mathbf{E} \to \mathbf{F}\) , \(v: \mathbf{F} \to \mathbf{G}\) two linear mappings.
\end{enumerate}

\begin{enumerate}
\def\labelenumi{(\alph{enumi})}
\tightlist
\item
  Show that there exists a decomposition of E as a direct sum of \(\operatorname{Ker}(u)\) and two subspaces M, N such that \(\operatorname{Ker}(v \circ u) = \mathbf{M} \oplus \operatorname{Ker}(u)\) and
\end{enumerate}

\[
\operatorname{Im} (v \circ u) = v (u (\mathbf {N})).
\]

\begin{enumerate}
\def\labelenumi{(\alph{enumi})}
\setcounter{enumi}{1}
\item
  \(\dim(\operatorname{Ker}(u)) \leqslant \dim(\operatorname{Ker}(v \circ u)) \leqslant \dim(\operatorname{Ker}(u)) + \dim(\operatorname{Ker}(v))\) ; if \(\operatorname{Ker}(u)\) and \(\operatorname{Ker}(v)\) are finite-dimensional, show that \(\dim(\operatorname{Ker}(v \circ u))\) can take every integer value satisfying the above inequalities.
\item
  The following equalities hold:
\end{enumerate}

\[
\begin{array}{l} \operatorname{rg} (u) = \operatorname{rg} (v \circ u) + \dim (\operatorname{Im} (u) \cap \operatorname{Ker} (v)) \\ \operatorname{rg} (v) = \operatorname{rg} (v \circ u) + \operatorname{codim} _ {\mathbb {F}} (\operatorname{Im} (u) + \operatorname{Ker} (v)). \end{array}
\]

If F is finite-dimensional, show that \(\operatorname{rg}(v \circ u)\) can take any integer value satisfying the inequality

\[
\sup (0, \operatorname{rg} (u) + \operatorname{rg} (v) - n) \leqslant \operatorname{rg} (v \circ u) \leqslant \inf (\operatorname{rg} (u), \operatorname{rg} (v)).
\]

\begin{enumerate}
\def\labelenumi{\arabic{enumi}.}
\setcounter{enumi}{13}
\item
  Let E, F, G be three vector spaces and \(u: \mathbf{E} \to \mathbf{F}\) , \(w: \mathbf{E} \to \mathbf{G}\) two linear mappings. Show that if \(u(0) \subset w(0)\) , there exists a linear mapping \(v\) of F into G such that \(w = v \circ u\) (cf.~§ 2, Exercise 8).
\item
  Let E, F be two vector spaces over a field K and \(u: E \rightarrow F\) a linear mapping.
\end{enumerate}

\begin{enumerate}
\def\labelenumi{(\alph{enumi})}
\item
  The following conditions are equivalent: (1) \(\operatorname{Ker}(u)\) is finite-dimensional; (2) there exists a linear mapping \(v\) of \(\mathbf{F}\) into \(\mathbf{E}\) such that \(v \circ u = 1_{\mathbf{E}} + w\) , where \(w\) is an endomorphism of \(\mathbf{E}\) of finite rank; (3) for every vector space \(\mathbf{G}\) over \(\mathbf{K}\) and every linear mapping \(f: \mathbf{G} \to \mathbf{E}\) , the relation \(\operatorname{rg}(u \circ f) < +\infty\) implies \(\operatorname{rg}(f) < +\infty\) .
\item
  The following conditions are equivalent: (1) Coker(u) is finite-dimensional; (2) there exists a linear mapping v of F into E such that \(u \circ v = 1_{F} + w\) , where w is an endomorphism of F of finite rank; (3) for every vector space G over K and every linear mapping \(g: F \to G\) , the relation \(\operatorname{rg}(g \circ u) < +\infty\) implies \(\operatorname{rg}(g) < +\infty\) .
\end{enumerate}

\begin{enumerate}
\def\labelenumi{\arabic{enumi}.}
\setcounter{enumi}{15}
\tightlist
\item
  Let E, F be two vector spaces over K and \(u: E \to F\) a linear mapping. u is said to be of finite index if \(\operatorname{Ker}(u)\) and \(\operatorname{Coker}(u)\) are finite-dimensional and the number \(d(u) = \dim(\operatorname{Ker}(u)) - \dim(\operatorname{Coker}(u))\) is then called the index of u.
\end{enumerate}

Let G be a third vector space over K and \(v: F \rightarrow G\) a linear mapping. Show that if two of the three linear mappings u, v, \(v \circ u\) are of finite index, so is the third and \(d(v \circ u) = d(u) + d(v)\) (use Exercise 13 (a)).

\begin{enumerate}
\def\labelenumi{\arabic{enumi}.}
\setcounter{enumi}{16}
\tightlist
\item
  Let E, F, G, H be four vector spaces over K and \(u: E \rightarrow F\) , \(v: F \rightarrow G\) , \(w: G \rightarrow H\) three linear mappings. Show that
\end{enumerate}

\[
\operatorname{rg} (v \circ u) + \operatorname{rg} (w \circ v) \leqslant \operatorname{rg} (v) + \operatorname{rg} (w \circ v \circ u).
\]

\begin{enumerate}
\def\labelenumi{\arabic{enumi}.}
\setcounter{enumi}{17}
\item
  Let E, F be two vector spaces over a field K and u, v two linear mappings of E into F. For there to exist an automorphism f of E and an automorphism g of F such that \(v = g \circ u \circ f\) , it is necessary and sufficient that \(\operatorname{rg}(u) = \operatorname{rg}(v)\) , \(\dim(\operatorname{Ker}(u)) = \dim(\operatorname{Ker}(v))\) and \(\dim(\operatorname{Coker}(u)) = \dim(\operatorname{Coker}(v))\) .
\item
  Let E be a vector space over a field K.
\end{enumerate}

\begin{enumerate}
\def\labelenumi{(\alph{enumi})}
\item
  Show that every mapping \(f\) of \(\mathbf{E}\) into \(\mathbf{E}\) which is permutable with every automorphism of \(\mathbf{E}\) is of the form \(x \mapsto \alpha x\) , where \(\alpha \in \mathbf{K}\) (show first that for all \(x \in \mathbf{E}\) there exists \(\rho(x) \in \mathbf{K}\) such that \(f(x) = \rho(x)x\) , using the fact that \(f\) commutes with every automorphism of \(\mathbf{E}\) leaving \(x\) invariant).
\item
  Let \(g\) be a mapping of \(\mathbf{E} \times \mathbf{E}\) into \(\mathbf{E}\) such that, for every automorphism \(u\) of \(\mathbf{E}\) , \(g(u(x), u(y)) = u(g(x, y))\) for all \(x, y\) in \(\mathbf{E}\) . Show that there exist two elements \(\alpha, \beta\) of \(\mathbf{K}\) such that \(g(x, y) = \alpha x + \beta y\) on the set of ordered pairs \((x, y)\) of linearly independent elements of \(\mathbf{E}\) ; moreover there exists a mapping \(\phi\) of \(\mathbf{K} \times \mathbf{K}\) into \(\mathbf{K}\) such that \(g(\lambda x, \mu x) = \phi(\lambda, \mu)x\) for all \(x \in \mathbf{E}\) (same method). If further \(g(u(x), u(y)) = u(g(x, y))\) for every endomorphism \(u\) of \(\mathbf{E}\) , then \(g(x, y) = \alpha x + \beta y\) for all \(x, y\) in \(\mathbf{E}\) . Generalize to mappings of \(\mathbf{E}^n\) into \(\mathbf{E}\) .
\end{enumerate}

\begin{enumerate}
\def\labelenumi{\arabic{enumi}.}
\setcounter{enumi}{19}
\tightlist
\item
  Let E be a vector space.
\end{enumerate}

\begin{enumerate}
\def\labelenumi{(\alph{enumi})}
\item
  Show that if M and N are two vector subspaces of E and M', N' the orthogonals of \(\mathbf{M}\) and \(\mathbf{N}\) respectively in \(\mathbf{E}^*\) , the orthogonal of \(\mathbf{M} \cap \mathbf{N}\) is \(\mathbf{M}' + \mathbf{N}'\) (cf.~§ 2, Exercise 9 (a)).
\item
  If E is infinite-dimensional, show that there exist hyperplanes \(H'\) of \(E^{*}\) such that the subspace of E orthogonal to \(H'\) reduces to 0 (consider a hyperplane containing the coordinate forms corresponding to a basis of E).
\item
  Show that if \(\mathbf{E}\) is infinite-dimensional there exists an infinite family \((\mathbf{V}_i)\) of subspaces of \(\mathbf{E}\) such that, if \(\mathbf{V}_i'\) is the subspace of \(\mathbf{E}^*\) orthogonal to \(\mathbf{V}_i\) , the subspace of \(\mathbf{E}^*\) orthogonal to \(\bigcap_{i} \mathbf{V}_i\) is distinct from \(\sum_{i} \mathbf{V}_i'\) .
\item
  Deduce from (b) that if E is infinite-dimensional, there exists a decomposition \(V' \oplus W'\) of \(E^{*}\) as a direct sum, \(W'\) being finite-dimensional, for which the sum \(V + W\) of subspaces V, W of E orthogonal respectively to \(V'\) and \(W'\) is distinct from E.
\item
  For a subspace of E to be finite-dimensional, it is necessary and sufficient that its orthogonal in \(E^{*}\) be of finite codimension.
\end{enumerate}

\begin{enumerate}
\def\labelenumi{\arabic{enumi}.}
\setcounter{enumi}{20}
\item
  With the notation of Theorem 8 of no. 6, let \(u\) denote the linear mapping \(x \mapsto (\langle x, x_i^* \rangle)\) of \(\mathbf{E}\) into \(\mathbf{K}_s^{\mathbf{I}}\) and \(y_0 = (\eta_i)\) . \(\mathbf{K}_d^{(\mathbf{I})}\) is identified with a subspace of the dual of \(\mathbf{K}_s^{\mathbf{I}}\) under the canonical mapping of \(\mathbf{K}_d^{(\mathbf{I})}\) into its bidual; then show that if \(\mathbf{N}'\) is the kernel of \(^t u\) , the condition of Theorem 3 (no. 5) expresses the fact that \(y_0\) is orthogonal to the intersection of \(\mathbf{N}'\) and \(\mathbf{K}_d^{(\mathbf{I})}\) . When \(\mathbf{I}\) is infinite and the system (21) (no. 5) is of finite rank, show that this intersection is distinct from \(\mathbf{N}'\) (note that \(\mathbf{N}'\) is then of finite codimension).
\item
  Let E and F be two vector spaces over a field K and u a linear mapping of E into F. If V is a vector subspace of E and \(V'\) the orthogonal of V in \(E^{*}\) , show that the dual of \(u(V)\) is isomorphic to \({}^{t}u(F^{*})/(V'\cap{}^{t}u(F^{*}))\) . If \(W'\) is a subspace of \(F^{*}\) such that \({}^{t}u(W')\) is finite-dimensional and W is the orthogonal of \(W'\) in F, \({}^{t}u(W')\) is isomorphic to the dual of the space \(u(E)/(W\cap u(E))\) .
\item
  Show that for a linear mapping \(u\) of a vector space \(E\) into a vector space \(F\) to be such that \(\mathrm{rg}(t u) = \mathrm{rg}(u)\) , it is necessary and sufficient that \(\mathrm{rg}(u)\) be finite (cf.~Exercise 3 (d)).
\item
  Let E be a vector space of finite dimension n \textgreater{} 1 over a commutative field K; show that, unless n = 2 and K is a field with two elements, there exists no isomorphism \(\phi\) of E onto \(E^{*}\) depending only on the vector space structure of E. (Note that if \(\phi\) is such an isomorphism, then necessarily \(\langle x, \phi(y) \rangle = \langle u(x), \phi(u(y)) \rangle\) for x, y in E and for every automorphism u of E (Set Theory, IV, § 1, no. 5).)
\item
  \begin{enumerate}
  \def\labelenumii{(\alph{enumii})}
  \tightlist
  \item
    Let K be a field and L, M two sets; there are canonical inclusions \((\mathbf{K}_{s}^{\mathbf{L}})^{(\mathbf{M})}\subset(\mathbf{K}_{s}^{(\mathbf{M})})^{\mathbf{L}}\subset\mathbf{K}_{s}^{\mathbf{L}\times\mathbf{M}}\) . Show that for two of these vector spaces to be equal, it is necessary and sufficient that one of the sets L, M be finite.
  \end{enumerate}
\end{enumerate}

\begin{enumerate}
\def\labelenumi{(\alph{enumi})}
\setcounter{enumi}{1}
\tightlist
\item
  Let \((\mathbf{E}_{\lambda})_{\lambda \in L}\) be a family of right vector spaces over \(K\) and \((\mathbf{F}_{\mu})_{\mu \in M}\) a family of left vector spaces over \(K\) . For the canonical mapping (26) (no. 7) to be bijective, it is necessary and sufficient that one of the vector spaces \(\prod_{\lambda \in L} \mathbf{E}_{\lambda}, \prod_{\mu \in M} F_{\mu}\) be finite-dimensional (use (a)).
\end{enumerate}

\begin{enumerate}
\def\labelenumi{\arabic{enumi}.}
\setcounter{enumi}{25}
\tightlist
\item
  \begin{enumerate}
  \def\labelenumii{(\alph{enumii})}
  \tightlist
  \item
    Let E, F be two left vector spaces over a field K; for the canonical homomorphism \(E^{*} \otimes_{K} F \to \operatorname{Hom}_{K}(E, F)\) to be bijective, it is necessary and sufficient that one of the spaces E, F be finite-dimensional.
  \end{enumerate}
\end{enumerate}

\begin{enumerate}
\def\labelenumi{(\alph{enumi})}
\setcounter{enumi}{1}
\tightlist
\item
  Let E be a right vector space and F a left vector space over a field K. For the canonical homomorphism \(E \otimes_{K} F \to \operatorname{Hom}_{K}(E^{*}, F)\) to be bijective, it is necessary and sufficient that E be finite-dimensional (use Exercise 3(d)).
\end{enumerate}

\begin{enumerate}
\def\labelenumi{\arabic{enumi}.}
\setcounter{enumi}{26}
\tightlist
\item
  \begin{enumerate}
  \def\labelenumii{(\alph{enumii})}
  \tightlist
  \item
    Under the hypotheses of Proposition 16 (i) (no. 7), for the canonical homomorphism (27) (no. 7) to be bijective, it is necessary and sufficient that E or F be finite-dimensional (note that the image under (27) (no. 7) of an element of \(\mathrm{Hom}_{\mathbb{K}}(\mathbf{E},\mathbf{G})\otimes_{\mathbb{L}}\mathbf{F}\) maps E to a subspace of \(\mathbf{G}\otimes_{\mathbb{L}}\mathbf{F}\) contained in a subspace of the form \(\mathbf{G}\otimes_{\mathbb{L}}\mathbf{F}'\) , where \(\mathbf{F}'\) is a finite-dimensional subspace of \(\mathbf{F}\) ).
  \end{enumerate}
\end{enumerate}

\begin{enumerate}
\def\labelenumi{(\alph{enumi})}
\setcounter{enumi}{1}
\tightlist
\item
  Under the hypotheses of Proposition 16 (ii) (no. 7), for the canonical homomorphism (28) (no. 7) to be bijective, it is necessary and sufficient that one of the ordered pairs \((\mathbf{E}_1, \mathbf{E}_2), (\mathbf{E}_1, \mathbf{F}_1), (\mathbf{E}_2, \mathbf{F}_2)\) consist of finite-dimensional vector spaces (use (a) and Exercise 7 of §4).
\end{enumerate}

\begin{enumerate}
\def\labelenumi{\arabic{enumi}.}
\setcounter{enumi}{27}
\item
  Let \(\rho: \mathbf{K} \to \mathbf{A}\) be an injective homomorphism of a field \(\mathbf{K}\) into a ring \(\mathbf{A}\) and let \(\mathbf{E}\) be a left vector space over \(\mathbf{K}\) . For the canonical mapping \((\mathbf{E}^*)_{(\mathbf{A})} \to (\mathbf{E}_{(\mathbf{A})})^*\) to be bijective, it is necessary and sufficient that \(\mathbf{E}\) be finite-dimensional or that \(\mathbf{A}\) , considered as a right vector \(\mathbf{K}\) -space, be finite-dimensional (use Exercise 25 ( \(b\) )). When one of these two cases holds, there exists a canonical bijection \((\mathbf{E}^{**})_{(\mathbf{A})} \to (\mathbf{E}_{(\mathbf{A})})^{**}\) .
\item
  Let A be an integral domain, K its field of fractions, E an A-module and T(E) its torsion module.
\end{enumerate}

\begin{enumerate}
\def\labelenumi{(\alph{enumi})}
\item
  Show that every linear form on \(\mathbf{E}\) is zero on \(\mathrm{T}(\mathrm{E})\) , so that the duals \(\mathbf{E}^*\) and \((\mathrm{E} / \mathrm{T}(\mathrm{E}))^*\) are canonically isomorphic.
\item
  Show that the canonical mapping \((\mathbf{E}^{*})_{(\mathbf{K})}\to (\mathbf{E}_{(\mathbf{K})})^{*}\) is injective and that it is bijective when \(\mathbf{E}\) is finitely generated.
\item
  Give an example of a torsion-free A-module E such that \(E^{*} = \{0\}\) and \(E_{(K)} \neq \{0\}\) .
\end{enumerate}

¶ 30. Let A be an integral domain and E, F two A-modules. Show that if x is a free element of E and y a free element of F, then \(x \otimes y \neq 0\) in \(E \otimes_{A} F\) . (Reduce it to the case where E and F are torsion-free, then to the case where E and F are finitely generated and use Exercise 29 (b) to show that there is an A-bilinear mapping f of \(E \times F\) into K such that \(f(x, y) \neq 0\) .)

*31. Let \(\mathbf{K}_0\) be a commutative field, A the polynomial ring \(\mathrm{K}_0[\mathrm{X},\mathrm{Y}]\) in two indeterminates over \(\mathbf{K}_0\) , which is an integral domain, and E the ideal \(\mathrm{AX} + \mathrm{AY}\) of A. Show that in the tensor product \(\mathbf{E} \otimes_{\mathbf{A}} \mathbf{E}\) the element \(\mathbf{X} \otimes \mathbf{Y} - \mathbf{Y} \otimes \mathbf{X}\) is \(\neq 0\) and that \(\mathrm{XY}(\mathbf{X} \otimes \mathbf{Y} - \mathbf{Y} \otimes \mathbf{X}) = 0\) (consider the A-bilinear mappings of \(\mathbf{E} \times \mathbf{E}\) into the quotient module \(\mathbf{A}/\mathbf{E}\) ).*

\begin{enumerate}
\def\labelenumi{\arabic{enumi}.}
\setcounter{enumi}{31}
\item
  Extend the results of no. 10 to the case where A is a non-commutative ring admitting a field of left fractions K (I, §9, Exercise 15). (Note that if \(\xi_{i}\) ( \(1 \leqslant i \leqslant n\) ) are elements of K, there exists an \(\alpha \neq 0\) in A such that all the \(\alpha\xi_{i}\) belong to A.)
\item
  Let A be an integral domain. An A-module E is called divisible if for all \(x \in E\) and all \(\alpha \neq 0\) in A, there exists \(y \in E\) such that \(\alpha y = x\) .
\end{enumerate}

\begin{enumerate}
\def\labelenumi{(\alph{enumi})}
\item
  Let E, F be two A-modules. Show that if E is divisible, so is \(E \otimes_{A} F\) . If E is divisible and F is torsion-free, \(\operatorname{Hom}_{\mathbf{A}}(\mathbf{E}, \mathbf{F})\) is torsion-free. If E is a torsion A-module and F is divisible, then \(E \otimes_{A} F = \{0\}\) . If E is a torsion A-module and F is torsion-free, then \(\operatorname{Hom}_{\mathbf{A}}(\mathbf{E}, \mathbf{F}) = \{0\}\) .
\item
  Show that every injective A-module (§ 2, Exercise 11) is divisible and conversely that every divisible torsion-free A-module is injective (use criterion (δ) of 2, Exercise 11).
\end{enumerate}

\begin{enumerate}
\def\labelenumi{\arabic{enumi}.}
\setcounter{enumi}{33}
\tightlist
\item
  Let A be an integral domain, P a projective A-module, \(P'\) a projective submodule of P and \(j: P' \to P\) the canonical injection. Show that for every torsion-free A-module E, the homomorphism \(j \otimes 1: P' \otimes_{A} E \to P \otimes_{A} E\) is injective (reduce it to the case where E is finitely generated and embed E in a free A-module).
\end{enumerate}

¶ 35. (a) Let K be a field and E a (K, K)-bimodule. Suppose that the dimensions of E, as a left and right vector space over K, are equal to the same finite number n.~Show that there exists a family \((e_{i})\) of n elements of E, which is a basis for each of the two vector space structures on E. (Note that if \((b_{j})_{1\leq j\leq m}\) is a family of m\textless n elements of E, which is free for each of the two vector space structures on E and V is the left vector subspace and W the right vector subspace generated by \((b_{j})\) , either \(V+W\neq E\) , or \(V\cap CW\) and \(W\cap CV\) are non-empty; in the latter case, if \(y\in V\cap CW\) , \(z\in W\cap CV\) , \(y+z\) forms with the \(b_{j}\) a family of \(m+1\) elements which is free for each of the two vector space structures on E.)

\begin{enumerate}
\def\labelenumi{(\alph{enumi})}
\setcounter{enumi}{1}
\item
  Let F be a subbimodule of E, whose two vector space structures over K have the same dimension p \textless{} n.~Show that there exists a family \((e_{i})_{1 \leqslant i \leqslant n}\) of elements of E which is a basis of E for each of its vector space structures and such that \((e_{i})_{1 \leqslant i \leqslant p}\) is a basis of F for each of its two vector space structures (same method).
\item
  Let \((b_{j})_{1\leqslant j\leqslant n - 1}\) be a family of \(n - 1\) elements of \(\mathbf{E}\) , which is free for each of the two vector space structures on \(\mathbf{E}\) . Let \(\mathbf{V}\) (resp. W) be the hyperplane generated by the \((b_{j})\) in \(\mathbf{E}\) considered as a left (resp. right) vector space.
\end{enumerate}

Show that if \(\mathbf{V} \subset \mathbf{W}\) (resp. \(\mathbf{W} \subset \mathbf{V}\) ), then necessarily \(\mathbf{V} = \mathbf{W}\) (note that if \(a \notin \mathbf{W}\) , the set of \(\lambda \in \mathbf{K}\) such that \(\lambda a \in \mathbf{W}\) is a left ideal).

*36. (a) Let \(\mathbf{K}_0\) be a commutative field and \(\mathbf{K} = \mathbf{K}_0(\mathbf{X})\) the field of rational functions in one indeterminate over \(\mathbf{K}_0\) (Chapter IV). A (K, K)-bimodule structure is defined on K as follows: the product \(t.u\) of an element \(u \in \mathbf{K}\) by a left operator \(t \in \mathbf{K}\) is the rational function \(t(\mathbf{X})u(\mathbf{X})\) ; the product \(u.t\) of \(u\) by a right operator \(t \in \mathbf{K}\) is the rational function \(u(\mathbf{X})t(\mathbf{X}^2)\) ; the left (resp. right) vector K-space structure thus defined on K is then 1-dimensional (resp. 2-dimensional). Deduce examples of (K, K)-bimodules such that dimensions of the two vector space structures of such a bimodule are arbitrary integers.

\begin{enumerate}
\def\labelenumi{(\alph{enumi})}
\setcounter{enumi}{1}
\tightlist
\item
  Deduce from (a) an example of a (K, K)-bimodule E whose two vector space structures have the same dimension and such that there exists a sub-bimodule F of E whose two vector space structures do not have the same dimension.*
\end{enumerate}

¶ 37. (a) Let E be a left vector space (finite-dimensional or otherwise) and \(A_{i}\) ( \(1 \leqslant i \leqslant n\) ) vector subspaces of E. Suppose that, for every sequence \((a_{i})_{1 \leqslant i \leqslant n}\) of points of E such that \(a_{i} \in A_{i}\) for all i, the vector subspace of E generated by the \(a_{i}\) is of dimension \(\leqslant m\) (m an integer \(\leqslant n\) ). Then prove that there exists a subspace W of E of dimension \(h \leqslant m\) , containing \(h + (n - m)\) of the subspaces \(A_{i}\) . (Argue by induction on n.~Prove first that (for fixed n) attention may be confined to the case where m \textless{} n and \(\dim(A_{i}) \leqslant m\) for all i; it may also be assumed that \(A_{i} \neq \{0\}\) for all i and that the \(A_{i}\) are not all

of dimension 1. Then argue (for fixed \(n\) ) by induction on \(d = \sum_{i=1}^{m} \dim(A_i)\) . If for example \(\dim(A_n) \geqslant 2\) , consider a 1-dimensional subspace \(B_n\) of \(A_n\) of dimension 1 and apply the induction hypothesis to \(A_1, \ldots, A_{n-1}, B_n\) ; conclude that there exists a number \(k\) such that \(1 \leqslant k \leqslant m\) and a \(k\) -dimensional subspace U of E containing \(k\) of the \(A_i\) , for example \(A_1, \ldots, A_k\) ; replacing, if need be, \(k\) by an integer \(k'\) such that \(1 \leqslant k' \leqslant k\) , show also that it can be assumed that there exist vectors \(b_i \in A_i\) for \(1 \leqslant i \leqslant k\) such that \(b_1, \ldots, b_k\) form a basis of U; then project onto a supplementary subspace of U in E and use the induction hypothesis.)

\begin{enumerate}
\def\labelenumi{(\alph{enumi})}
\setcounter{enumi}{1}
\tightlist
\item
  Let K be a finite field and let E be an N-dimensional vector space over K. Let \(U_{i}\) ( \(1 \leqslant i \leqslant n\) ) be vector subspaces of E such that, for every subset H of \([1, n]\) , the intersection of the \(U_{i}\) such that \(i \in H\) has dimension \(\leqslant N - \text{Card}(H)\) . Show that the union of the \(U_{i}\) cannot be equal to E (argue by duality using (a)).
\end{enumerate}

\begin{enumerate}
\def\labelenumi{\arabic{enumi}.}
\setcounter{enumi}{37}
\item
  Let K be a commutative field of characteristic 0 and E a finite-dimensional vector space over K. Let \(p_{1}, \ldots, p_{m}\) be projectors of E onto vector subspaces of E such that \(1_{E} = p_{1} + p_{2} + \cdots + p_{m}\) . Show that E is the direct sum of the \(p_i(\mathbf{E})\) and that the \(p_i\) are pairwise permutable (take the traces of the two sides).
\item
  Let K be a field, L a subfield of K and \((\mathrm{K}_{\alpha})_{\alpha\in I}\) a left directed family of subfields of K such that \(L=\bigcap_{\alpha\in I}K_{\alpha}\) . Let E be a vector space over K and \((a_{i})_{1\leqslant i\leqslant m}\) a finite family of elements of E; show that if the family \((a_{i})\) is free over L, there exists an index \(\alpha\) such that \((a_{i})\) is free over \(K_{\alpha}\) . (Argue by induction on m-r, where r is the rank of the family \((a_{i})\) over K.)
\end{enumerate}

\BourbakiExercisesSection

\begin{enumerate}
\def\labelenumi{\arabic{enumi}.}
\tightlist
\item
  Let K be a field, \(K'\) a subfield of K and V a non-zero right vector space over K.
\end{enumerate}

\begin{enumerate}
\def\labelenumi{(\alph{enumi})}
\item
  Let \(\mathbf{V}'\) be a \(\mathbf{K}'\) -structure on \(\mathbf{V}\) . Show that \(\mathbf{K}'\) is equal to the set of \(\mu \in \mathbf{K}\) such that, for all \(x \in \mathbf{V}'\) , \(x\mu \in \mathbf{V}'\) .
\item
  Let \(\Gamma'\) be the group of automorphisms of K leaving invariant the elements of K'. Show that if there exists no element of K not belonging to K' and invariant under all the automorphisms \(\sigma \in \Gamma'\) , there exists no element of V not belonging to V' and invariant under every bijective dimorphism of V (relative to an automorphism of K) which leaves invariant all the elements of V'.
\item
  Conversely, let E be a subset of V and G the group of bijective dimorphisms of V leaving invariant all the elements of E; for all \(u \in G\) , let \(\sigma_{u}\) be the corresponding automorphism of K and let \(\Gamma\) be the group of automorphisms of K the image of G under the homomorphism \(u \mapsto \sigma_{u}\) . Suppose that G contains no automorphism of V distinct from the identity automorphism and that there exists no element of V not belonging to E and invariant under all the dimorphisms \(u \in G\) . Show that if \(K'\) is the subfield of elements of K invariant under all the automorphism \(\sigma \in \Gamma\) , E is a \(K'\) -structure on V. (Prove first that E is an additive subgroup of V containing a basis of V; let \(K''\) be the set of elements \(\mu \in K\) such that the relation \(x \in E\) implies \(x\mu \in E\) . Show that the elements of \(K''\) are invariant under all \(\sigma \in \Gamma\) and that \(K''\) is a subfield of K; deduce that E is a \(K''\) -structure on V.)
\end{enumerate}

\begin{enumerate}
\def\labelenumi{\arabic{enumi}.}
\setcounter{enumi}{1}
\tightlist
\item
  Let K be a field, \(K'\) a subfield of K, V a right vector space over K, \(V'\) a \(K'\) -structure on V and \(R'\) the canonical image of the dual \(V'*\) in the dual \(V*\) of V (no. 4). For \(R'\) to be a generating system of \(V*\) , it is necessary and sufficient that V be finite-dimensional or that K be a finite-dimensional right vector \(K'\) -space (use Exercise 25 of §7).
\end{enumerate}

¶ 3. Let K be a field, L a subfield of K and let \(K_{L}\) denote the field K considered as a left vector space over L; \(E = \text{End}_{L}(K_{L})\) has canonically a (K, K)-bimodule structure. The dual \((K_{L})^{*}\) is contained in E and is a (K, L)-sub-bimodule of E. When L is contained in the centre of K, the two vector L-space structures on E, obtained by restricting to L the field of scalars of the two vector K-space structures on E, are identical.

Suppose henceforth that \(\mathbf{L}\) is contained in the centre of \(\mathbf{K}\) and that the dimension of \(\mathbf{K}_{\mathbf{L}}\) is finite and equal to \(n\) .

\begin{enumerate}
\def\labelenumi{(\alph{enumi})}
\item
  Show that \((\mathrm{K}_{\mathrm{L}})^{*}\) is 1-dimensional under its left vector space structure over K (calculate the dimension of \((\mathrm{K}_{\mathrm{L}})^{*}\) over L in two ways).
\item
  Show that \(\mathbf{E}\) is an \(n\) -dimensional left vector space over \(\mathbf{K}\) (same method).
\item
  Let F be a finite-dimensional left vector space over K and \(F_{[L]}\) the corresponding vector space over L. If \(u_{0} \neq 0\) is a linear form on \(K_{L}\) , show that the mapping \(x' \mapsto u_{0} \circ x'\) of the dual \(F^{*}\) of F (considered as a vector space over L) onto the dual \((\mathrm{F}_{[\mathrm{L}]})^{*}\) of \(F_{[L]}\) is bijective (use (a)).
\end{enumerate}

¶ 4. Let K be a field of finite rank over its centre Z and L a subfield of K containing Z.

\begin{enumerate}
\def\labelenumi{(\alph{enumi})}
\item
  Show that the left and right vector space structures of K with respect to L have the same dimension.
\item
  Show that properties (a), (b), (c) of Exercise 3 also hold in this case (if \(u_0\) is a linear form \(\neq 0\) on \(\mathbf{L}_{\mathbf{Z}}\) and \(v_0\) a linear form \(\neq 0\) on \(\mathbf{K}_{\mathbf{L}}\) , note that \(\xi. (u_0 \circ v_0) = u_0 \circ (\xi v_0)\) for \(\xi \in \mathbf{K}\) , that \(\xi. (u_0 \circ v_0)\) runs through the dual \((\mathbf{K}_{\mathbf{Z}})^*\) when \(\xi\) runs through \(\mathbf{K}\) and use Exercise 3 (c)).
\end{enumerate}

\begin{enumerate}
\def\labelenumi{\arabic{enumi}.}
\setcounter{enumi}{4}
\tightlist
\item
  Let \(K_{0}\) be a subfield of a field K such that K is of dimension 2 as a right vector space over \(K_{0}\) . Let E be a left vector space over K.
\end{enumerate}

\begin{enumerate}
\def\labelenumi{(\alph{enumi})}
\item
  Let \(E_{0}\) be a subset of E which is a vector space over \(K_{0}\) and let V be the largest vector sub-K-space of \(E_{0}\) ; if \(W_{0}\) is a vector sub- \(K_{0}\) -space of \(E_{0}\) supplementary to V, show that the vector sub-K-space W of E generated by \(W_{0}\) is such that \(V \cap W = \{0\}\) (note that if an element \(\mu \in K\) does not belong to \(K_{0}\) , the relation \(\mu x \in E_{0}\) for an \(x \in E_{0}\) implies \(x \in V\) ).
\item
  Let \(E_{0}^{\prime}\) be another vector sub- \(K_{0}\) -space of E and \(V^{\prime}\) the largest vector sub-K-space of \(E_{0}^{\prime}\) . For there to exist a K-automorphism of E mapping \(E_{0}\) to \(E_{0}^{\prime}\) , it is necessary and sufficient that V and \(V^{\prime}\) have the same dimension with respect to K, that the codimensions of V in \(E_{0}\) and \(V^{\prime}\) in \(E_{0}^{\prime}\) (with respect to \(K_{0}\) ) be equal and that the codimensions of \(E_{0}\) and \(E_{0}^{\prime}\) in E (with respect to \(K_{0}\) ) be equal (use (a)).
\end{enumerate}

\BourbakiExercisesSection

*1. In an affine space \(\mathbf{E}\) over a field \(\mathbf{K}\) , a quadruple \((a, b, c, d)\) of points of \(\mathbf{E}\) is called a parallelogram if \(b - a = c - d\) , in which case \((a, d, c, b)\) is also a parallelogram. Show that if \(\mathbf{K}\) is of characteristic \(\neq 2\) the midpoints of the ordered pairs \((a, c)\) and \((b, d)\) are then equal; what can be said when \(\mathbf{K}\) is of characteristic \(2?_{*}\) .

*2. Let E be an affine space over a field of characteristic \(\neq 2\) and \(a, b, c, d\) any four points of E. Show that if \(x, y, z, t\) are the respective midpoints of the ordered pairs \((a, b), (b, c), (c, d), (d, a)\) , the quadruple \((x, y, z, t)\) is a parallelogram (Exercise 1).*

*3. Let K be a commutative field of characteristic \(\neq 2\) , E an affine plane over K and \(a, b, c, d\) four points of E any three of which do not lie on a straight line. Let \(\mathbf{D}_{xy}\) denote the line passing through two distinct points \(x, y\) of E. Suppose that the lines \(\mathbf{D}_{ab}\) and \(\mathbf{D}_{cd}\) have a common point \(e\) and that \(\mathbf{D}_{ad}\) and \(\mathbf{D}_{bc}\) have a common point \(f\) . Show that the midpoints of the three ordered pairs \((a, c), (b, d), (e, f)\) lie on a straight line. What becomes of this property when \(\mathbf{D}_{ab}\) and \(\mathbf{D}_{cd}\) are parallel or when \(\mathbf{D}_{ad}\) and \(\mathbf{D}_{bc}\) are parallel? Consider the case where K is a field with 3 elements.*

*4. Let K be a field whose characteristic is different from 2 and 3, E an affine space over K, \(a, b, c\) three points of E not on a straight line and \(a', b', c'\) the respective midpoints of the ordered pairs \((b, c)\) , \((c, a)\) and \((a, b)\) . Show that (in the notation of Exercise 3) the lines \(\mathbf{D}_{aa'}\) , \(\mathbf{D}_{bb'}\) and \(\mathbf{D}_{cc'}\) pass through the barycentre of the three points \(a, b, c\) . What becomes of this property when K is of characteristic 2 or of characteristic 3? Generalize to a system of \(n\) affinely independent points.*

\begin{enumerate}
\def\labelenumi{\arabic{enumi}.}
\setcounter{enumi}{4}
\item
  For a non-empty subset V of an affine space E over a field K with at least three elements to be a linear variety, it is necessary and sufficient that for every ordered pair \((x,y)\) of distinct points of V the line \(D_{xy}\) passing through x and y be entirely contained in V. If K has two elements, for V to be a linear variety, it is necessary and sufficient that the barycentre of any three points of V belong to V.
\item
  \begin{enumerate}
  \def\labelenumii{(\alph{enumii})}
  \tightlist
  \item
    Let E be an affine space of dimension \(\geqslant2\) over a field K. For an affine mapping u of E into itself to transform every line of E into a parallel line, it is necessary and sufficient that the linear mapping v associated with u be a homothety \(t\mapsto\gamma t\) of ratio \(\gamma\neq0\) belonging to the centre of K. If \(\gamma=1\) , u is a translation; show that if \(\gamma\neq1\) , there exists one and only one point \(a\in E\) such that \(u(a)=a\) . If a is taken as origin of E, u is then identified with a central homothety for the vector space structure thus determined on E; u is called a central homothety of the affine space E of centre a and ratio \(\gamma\) .
  \end{enumerate}
\end{enumerate}

\begin{enumerate}
\def\labelenumi{(\alph{enumi})}
\setcounter{enumi}{1}
\item
  Let \(u_1, u_2\) be two affine mappings of \(\mathbf{E}\) into \(\mathbf{E}\) each of which is either a translation or a central homothety. Show that \(u_1 \circ u_2\) is a translation or a central homothety of \(\mathbf{E}\) ; if \(u_1, u_2\) and \(u_1 \circ u_2\) are all three central homotheties show that their centres lie on a straight line. What can be said when \(u_1\) and \(u_2\) are central homotheties and \(u_1 \circ u_2\) is a translation?
\item
  Show that the set of translations and central homotheties is a normal subgroup H of the affine group of E and that H/T is isomorphic to the multiplicative group of the centre of K: show that H can only be commutative if H = T, in other words if the centre of K has only two elements.
\end{enumerate}

¶ 7. Let E (resp. E') be an affine space of finite dimension \(n \geqslant 2\) over field K with at least 3 elements (resp. over a field K') and let u be an injective mapping of E into E' mapping any three points in a straight line in E to three points in a straight line and such that the affine linear variety generated by \(u(E)\) in E' is equal to E'.

\begin{enumerate}
\def\labelenumi{(\alph{enumi})}
\item
  Show that \(u\) maps every system of affinely independent points of \(E\) to a system of affinely independent points (use Exercise 5).
\item
  Let \(D_{1}, D_{2}\) be two parallel lines in E and \(D_{1}^{\prime}, D_{2}^{\prime}\) the lines of \(E^{\prime}\) containing respectively \(u(\mathbf{D}_{1})\) and \(u(\mathbf{D}_{2})\) . Show that \(D_{1}^{\prime}\) and \(D_{2}^{\prime}\) are in the same plane; if further u is surjective, show that \(D_{1}^{\prime}\) and \(D_{2}^{\prime}\) are parallel (in the contrary case, show that there would be 3 points not in a straight line in E whose images under u would be in a straight line) (cf.~Exercise 17).
\item
  Suppose henceforth that if \(D_{1}\) and \(D_{2}\) are parallel lines in E, the lines of \(E'\) containing respectively \(u(\mathbf{D}_{1})\) and \(u(\mathbf{D}_{2})\) are parallel. Show that if an origin a in E and the origin \(a' = u(a)\) in \(E'\) are taken, there exists an isomorphism \(\sigma\) of K onto a subfield \(K_{1}\) of \(K'\) such that if E is considered as a vector space over K and \(E'\) as a vector space over \(K_{1}, u\) is an injective semi-linear mapping (relative to \(\sigma\) ) of E into \(E'\) (§ 1, no. 13). (Consider first the case n = 2; given a basis \((e_{1}, e_{2})\) of E, show that for any two elements \(\alpha, \beta\) of K, the points \((\alpha + \beta)e_{1}\) and \((\alpha\beta)e_{1}\) of E can be constructed starting with the points 0, \(e_{1}, e_{2}, \alpha e_{1}, \beta e_{1}\) by constructing parallels to given lines and intersections of given lines; deduce that \(u(\lambda e_{1}) = \lambda^{\sigma}u(e_{1})\) , where \(\sigma\) is an isomorphism of K onto a subfield of \(K'\) , then show that also \(u(\lambda e_{2}) = \lambda^{\sigma}u(e_{2})\) by considering the line joining the points \(\lambda e_{1}\) and \(\lambda e_{2}\) . Pass finally from here to the case where n is arbitrary, arguing by induction on n.) If u is bijective, show that \(K_{1} = K'\) .
\item
  Extend the result of (c) to the case where K is a field with 2 elements assuming also that u maps every system of affinely independent points of E to a system of affinely independent points of E'.
\end{enumerate}

\begin{enumerate}
\def\labelenumi{\arabic{enumi}.}
\setcounter{enumi}{7}
\tightlist
\item
  Let E be a left affine space over a field K and T its translation space.
\end{enumerate}

\begin{enumerate}
\def\labelenumi{(\alph{enumi})}
\tightlist
\item
  For a mapping \(f\) of \(\mathbf{E}\) into a left vector space \(\mathbf{L}\) over \(\mathbf{K}\) to be affine, it is necessary and sufficient that
\end{enumerate}

\[
\begin{array}{r l} f (\mathbf {t} + x) - f (x) & = f (\mathbf {t} + y) - f (y) \\ f (\lambda \mathbf {t} + x) - f (x) & = \lambda f (\mathbf {t} + x) - f (x)) \end{array}
\]

for all \(\lambda\in\mathbf{K},\mathbf{t}\in\mathbf{T},x,y\) in E. Let \(x\mapsto[x]\) be the canonical injection of E into the vector space \(\mathrm{K}_{s}^{(\mathbb{E})}\) of formal linear combinations of elements of E and let N be the subspace of \(\mathrm{K}_{s}^{(\mathbb{E})}\) generated by the elements

\[
[ \mathbf {t} + x ] - [ x ] - [ \mathbf {t} + y ] + [ y ]
\]

\[
[ \lambda \mathbf {t} + x ] - [ x ] - \lambda [ \mathbf {t} + x ] + \lambda [ x ]
\]

for \(\lambda \in \mathbf{K}\) , \(\mathbf{t} \in \mathbf{T}\) , \(x, y\) in \(\mathbf{E}\) ; finally, let \(\mathbf{V}\) be the quotient space of \(\mathbf{K}_s^{(\mathbf{E})}\) by \(\mathbf{N}\) and \(\phi\) the mapping of \(\mathbf{E}\) into \(\mathbf{V}\) which associates with every \(x \in \mathbf{E}\) the class of \([x]\) modulo \(\mathbf{N}\) . Then \(\phi\) is an affine mapping and for every affine mapping \(f\) of \(\mathbf{E}\) into a left vector space \(\mathbf{L}\) over \(\mathbf{K}\) , there exists one and only one linear mapping \(g\) of \(\mathbf{V}\) into \(\mathbf{L}\) such that \(f = g \circ \phi\) .

\begin{enumerate}
\def\labelenumi{(\alph{enumi})}
\setcounter{enumi}{1}
\tightlist
\item
  Let \(\phi_0: \mathbf{T} \to \mathbf{V}\) be the linear mapping associated with \(\phi\) , so that \(\phi(\mathbf{t} + x) - \phi(x) = \phi_0(\mathbf{t})\) . Show that \(\phi\) (and therefore also \(\phi_0\) ) is injective (consider the mapping \(x \mapsto x - a\) of \(\mathbf{E}\) into \(\mathbf{T}\) for some \(a \in \mathbf{E}\) ); for every family \((\lambda_t)_{t \in I}\) of elements of \(\mathbf{K}\) with finite support and every family \((x_t)_{t \in I}\) of elements of \(\mathbf{E}\) ,
\end{enumerate}

\[
\sum_ {i} \lambda_ {i} \phi (x _ {i}) = \phi_ {0} \left(\sum_ {i} \lambda_ {i} x _ {i}\right) \quad \text { if } \quad \sum_ {i} \lambda_ {i} = 0
\]

\[
\sum_ {i} \lambda_ {i} \phi (x _ {i}) = \mu \phi_ {0} \left(\sum_ {i} \mu^ {- 1} \lambda_ {i} x _ {i}\right) \quad \text { if } \quad \sum_ {i} \lambda_ {i} = \mu \neq 0.
\]

Deduce that \(\phi_{0}(T)\) is a hyperplane passing through 0 in V and \(\phi(E)\) a hyperplane parallel to \(\phi_{0}(T)\) .

\begin{enumerate}
\def\labelenumi{\arabic{enumi}.}
\setcounter{enumi}{8}
\tightlist
\item
  Let \(\mathbf{K}\) be a finite field with \(q\) elements and \(\mathbf{V}\) an \(n\) -dimensional vector space over \(\mathbf{K}\) .
\end{enumerate}

\begin{enumerate}
\def\labelenumi{(\alph{enumi})}
\tightlist
\item
  Show that the set of sequences \((x_{1}, x_{2}, \ldots, x_{m})\) of \(m \leqslant n\) vectors of \(\mathbf{V}\) forming a free system has cardinal equal to
\end{enumerate}

\[
(q ^ {n} - 1) (q ^ {n} - q) \dots (q ^ {n} - q ^ {m - 1})
\]

(argue by induction on \(n\) ).

\begin{enumerate}
\def\labelenumi{(\alph{enumi})}
\setcounter{enumi}{1}
\tightlist
\item
  Deduce from (a) that the cardinal of the set of \(m\) -dimensional linear varieties in an \(n\) -dimensional projective space over \(K\) is equal to
\end{enumerate}

\[
\frac {(q ^ {n + 1} - 1) (q ^ {n + 1} - q) \dots (q ^ {n + 1} - q ^ {m})}{(q ^ {m + 1} - 1) (q ^ {m + 1} - q) \dots (q ^ {m + 1} - q ^ {m})}.
\]

\begin{enumerate}
\def\labelenumi{\arabic{enumi}.}
\setcounter{enumi}{9}
\item
  In an \(n\) -dimensional projective space \(\mathbf{P}(\mathbf{V})\) over a field \(\mathbf{K}\) , a projective basis is a set \(\mathbf{S}\) of \(n + 2\) points any two of which form a projectively free system. If \(\mathbf{S} = (a_i)_{0 \leqslant i \leqslant n + 1}\) and \(\mathbf{S}' = (a_i')_{0 \leqslant i \leqslant n + 1}\) are any two projective bases of \(\mathbf{P}(\mathbf{V})\) , show that there exists a transformation \(f \in \mathbf{PGL}(\mathbf{V})\) such that \(f(a_i) = a_i'\) for \(0 \leqslant i \leqslant n + 1\) . For this transformation to be unique, it is necessary and sufficient that \(\mathbf{K}\) be commutative. (Reduce it to the case where \(a_i' = a_i\) for all \(i\) and note that it is always possible to write (in the notation of no. 6) \(a_i = \pi(b_i)\) , where \((b_i)_{1 \leqslant i \leqslant n + 1}\) is a basis of \(\mathbf{V}\) and \(b_0 = b_1 + b_2 + \cdots + b_{n+1}\) .) *Give an example where \(\mathbf{K}\) is a field of quaternions and there exists an infinite set \(\mathbf{T}\) of points of \(\mathbf{P}(\mathbf{V})\) any \(n + 1\) of which form a projectively free system and a transformation \(f \in \mathbf{PGL}(\mathbf{V})\) , distinct from the identity, leaving invariant all the points of \(\mathbf{T}\) .*
\item
  Let V be a 2-dimensional vector space over a field K and a, b, c, d four distinct points of the projective line \(\mathbf{P}(\mathbf{V})\) . The cross-ratio of the quadruple \((a, b, c, d)\) , denoted by \(\begin{bmatrix} a & b \\ d & c \end{bmatrix}\) , is the set of elements \(\xi \in \mathbf{K}\) such that there exist two vectors \(u, v\) in \(\mathbf{V}\) for which (in the notation of no. 6) \(a = \pi(u)\) , \(b = \pi(v)\) , \(c = \pi(u + v)\) , \(d = \pi(u + \xi v)\) . This definition extends immediately to any quadruple of distinct points of a set with a projective line structure (no. 11).
\end{enumerate}

\begin{enumerate}
\def\labelenumi{(\alph{enumi})}
\item
  Show that \(\begin{bmatrix} a & b \\ d & c \end{bmatrix}\) is the set of conjugates of an element \(\neq 1\) of the multiplicative group \(\mathbf{K}^*\) and conversely that, when \(a, b, c\) , are distinct points of \(\mathbf{P}(V)\) and \(\rho\) the set of conjugates of an element \(\neq 1\) of \(\mathbf{K}^*\) , there exists a point \(d \in \mathbf{P}(V)\) such that \(\begin{bmatrix} a & b \\ d & c \end{bmatrix} = \rho\) . For \(d\) to be unique, it is necessary and sufficient that \(\rho\) consist of a single point.
\item
  Show that
\end{enumerate}

\[
\left[ \begin{array}{c c} a & b \\ c & d \end{array} \right] = \left[ \begin{array}{c c} b & a \\ d & c \end{array} \right] = \left[ \begin{array}{c c} a & b \\ d & c \end{array} \right] ^ {- 1}
\]

and

\[
\left[ \begin{array}{c c} d & a \\ c & b \end{array} \right] = 1 - \left[ \begin{array}{c c} a & b \\ d & c \end{array} \right]
\]

(denoting by \(\rho^{-1}\) (resp. \(1 - \rho\) ) the set of conjugates \(\lambda \xi^{-1} \lambda^{-1} = (\lambda \xi \lambda^{-1})^{-1}\) (resp. \(1 - \lambda \xi \lambda^{-1} = \lambda (1 - \xi) \lambda^{-1}\) ), where \(\xi\) is an element of \(\rho\) ).

\begin{enumerate}
\def\labelenumi{(\alph{enumi})}
\setcounter{enumi}{2}
\tightlist
\item
  Let \((a, b, c, d)\) , \((a', b', c', d')\) be two quadruples of distinct points of \(\mathbf{P}(\mathbf{V})\) . For there to exist a bijective semi-linear mapping of V onto itself such that the bijective mapping f of \(\mathbf{P}(\mathbf{V})\) onto itself, obtained when passing to the quotients, satisfies the conditions \(f(a) = a'\) , \(f(b) = b'\) , \(f(c) = c'\) , \(f(d) = d'\) , it is necessary and sufficient that there exist an automorphism \(\sigma\) of K such that
\end{enumerate}

\[
\left[ \begin{array}{c c} a ^ {\prime} & b ^ {\prime} \\ d ^ {\prime} & c ^ {\prime} \end{array} \right] = \left[ \begin{array}{c c} a & b \\ d & c \end{array} \right] ^ {\sigma}.
\]

For there to exist a transformation f of the projective group PGL(V) satisfying the above conditions, it is necessary and sufficient that

\[
\left[ \begin{array}{c c} a ^ {\prime} & b ^ {\prime} \\ d ^ {\prime} & c ^ {\prime} \end{array} \right] = \left[ \begin{array}{c c} a & b \\ d & c \end{array} \right].
\]

\begin{enumerate}
\def\labelenumi{\arabic{enumi}.}
\setcounter{enumi}{11}
\tightlist
\item
  Let \(\mathbf{P}(\mathbf{V})\) be a (left) projective space of finite dimension \(n\) over a field \(\mathbf{K}\) . Show that there exists on the set of projective hyperplanes of \(\mathbf{P}(\mathbf{V})\) an \(n\) -dimensional (right) projective space structure over \(\mathbf{K}\) canonically isomorphic to that of \(\mathbf{P}(\mathbf{V}^*)\) ( \(\mathbf{V}^*\) being the dual of \(\mathbf{V}\) ). If \(\mathbf{M}\) is a linear variety of dimension \(r < n\) in \(\mathbf{P}(\mathbf{V})\) , derive an \((n - r - 1)\) -dimensional projective space structure on the set of projective hyperplanes containing \(\mathbf{M}\) . In particular if \(\mathbf{M}\) is of dimension \(n - 2\) , the cross-ratio \(\begin{bmatrix} \mathrm{H}_1 & \mathrm{H}_2 \\ \mathrm{H}_4 & \mathrm{H}_3 \end{bmatrix}\) of a quadruple \((\mathrm{H}_1, \mathrm{H}_2, \mathrm{H}_3, \mathrm{H}_4)\) of distinct hyperplanes containing M can be defined. Show that if \(\mathbf{D} \subset \mathbf{P}(\mathbf{V})\) is a line not meeting M and \(a_i\) is the intersection of D with \(\mathrm{H}_i\) ( \(1 \leqslant i \leqslant 4\) ), then
\end{enumerate}

\[
\left[ \begin{array}{c c} a _ {1} & a _ {2} \\ a _ {4} & a _ {3} \end{array} \right] = \left[ \begin{array}{c c} \mathrm{H} _ {1} & \mathrm{H} _ {2} \\ \mathrm{H} _ {4} & \mathrm{H} _ {3} \end{array} \right].
\]

*13. In a projective plane \(\mathbf{P}(\mathbf{V})\) over a field \(\mathbf{K}\) of characteristic \(\neq 2\) , let \(a, b, c, d\) be four points forming a projective basis (Exercise 10); denoting by \(\mathbf{D}_{xy}\) the line passing through two distinct points \(x, y\) of \(\mathbf{P}(\mathbf{V})\) , let \(e, f, g\) be the points of intersection of the lines \(\mathbf{D}_{ab}\) and \(\mathbf{D}_{cd}\) , \(\mathbf{D}_{ac}\) and \(\mathbf{D}_{bd}\) , \(\mathbf{D}_{ad}\) and \(\mathbf{D}_{bc}\) respectively; let \(h\) be the point of intersection of \(\mathbf{D}_{bc}\) and \(\mathbf{D}_{ef}\) ; show that \(\begin{bmatrix} b & c \\ h & g \end{bmatrix} = \{-1\}\) (``theorem of the complete quadrilateral''; reduce it to the case where \(\mathbf{D}_{ad}\) is the line at infinity of an affine plane). What is the corresponding result when \(\mathbf{K}\) is of characteristic \(2?_{*}\)

\begin{enumerate}
\def\labelenumi{\arabic{enumi}.}
\setcounter{enumi}{13}
\item
  In a projective plane \(\mathbf{P}(\mathrm{V})\) over a field \(\mathbf{K}\) with at least three elements, let \(\mathbf{D}, \mathbf{D}'\) be two distinct lines. In order that, for any distinct points \(a, b, c\) of \(\mathbf{D}, a', b', c'\) of \(\mathbf{D}'\) , the intersection points \(r\) of \(\mathbf{D}_{ab'}\) and \(\mathbf{D}_{ba'}\) , \(q\) of \(\mathbf{D}_{ac'}\) and \(\mathbf{D}_{ca'}\) , \(p\) of \(\mathbf{D}_{bc'}\) and \(\mathbf{D}_{cb'}\) lie on a straight line, it is necessary and sufficient that \(\mathbf{K}\) be commutative (``Pappus's theorem''; reduce it to the case where \(q\) and \(r\) are on the line at infinity of an affine plane). Apply this theorem to the projective space of lines of \(\mathbf{P}(\mathrm{V})\) (Exercise 12).
\item
  In a projective plane over a field K with at least three elements, let \(s, a, b, c, a', b', c'\) be seven distinct points such that \(\{s, a, b, c\}\) and \(\{s, a', b', c'\}\) are projective bases (Exercise 10) and the lines \(\mathbf{D}_{sa}, \mathbf{D}_{sb}, \mathbf{D}_{sc}\) pass respectively through \(a', b', c'\) . Show that the intersection points \(r\) of \(\mathbf{D}_{ab}\) and \(\mathbf{D}_{a'b'}\) , \(p\) of \(\mathbf{D}_{bc}\) and \(\mathbf{D}_{b'c'}\) , \(q\) of \(\mathbf{D}_{ca}\) and \(\mathbf{D}_{c'a'}\) lie on a straight line (``Desargues's theorem''; method analogous to that of Exercise 14).
\item
  \begin{enumerate}
  \def\labelenumii{(\alph{enumii})}
  \tightlist
  \item
    Let \(\mathbf{E} = \mathbf{P}(\mathbf{V})\) and \(\mathbf{E}' = \mathbf{P}(\mathbf{V}')\) be two projective spaces of the same dimension \(n \geqslant 2\) over two fields \(\mathbf{K}\) , \(\mathbf{K}'\) respectively and let \(u\) be a bijective mapping of \(\mathbf{E}\) onto \(\mathbf{E}'\) , mapping any three points on a straight line to three points on a straight line. Show that there exist an isomorphism \(\sigma\) of \(\mathbf{K}\) onto \(\mathbf{K}'\) and a bijective semi-linear mapping \(v\) of \(\mathbf{V}\) onto \(\mathbf{V}'\) (relative to \(\sigma\) ) such that \(u\) is the mapping obtained from \(v\) when passing to the quotients (``fundamental theorem of projective geometry''; use Exercise 7). Suppose also that \(\mathbf{V}' = \mathbf{V}\) and \(\mathbf{K}\) is commutative; for \(u\) to be a projective mapping, it is necessary and sufficient that also \(\begin{bmatrix} u(a) & u(b) \\ u(d) & u(c) \end{bmatrix} = \begin{bmatrix} a & b \\ d & c \end{bmatrix}\) for every quadruple \((a, b, c, d)\) of distinct points on a straight line in \(\mathbf{P}(\mathbf{V})\) .
  \end{enumerate}
\end{enumerate}

\begin{enumerate}
\def\labelenumi{(\alph{enumi})}
\setcounter{enumi}{1}
\tightlist
\item
  Let \(p\) be an integer such that \(1 \leqslant p \leqslant n - 1\) . Show that the first conclusion of (a) holds when the image under \(u\) of every \(p\) -dimensional projective linear variety is contained in a \(p\) -dimensional projective linear variety.
\end{enumerate}

\begin{enumerate}
\def\labelenumi{\arabic{enumi}.}
\setcounter{enumi}{16}
\tightlist
\item
  Let V be a vector space of finite dimension n over a field K, \((e_{i})_{1\leq i\leq n}\) a basis of V, \(K'\) a subfield of K and \(V'\) the n-dimensional vector space over \(K'\) generated by the \(e_{i}\) . Give an example of an injective mapping of \(V'\) into V, mapping any three points on a straight line of the affine space \(V'\) to three points on a straight line in the affine space V, but not necessarily mapping two parallel lines to sets contained in two parallel lines. (Embed V in the projective space E canonically associated with it and consider a projective transformation u of E into itself such that the inverse image under u of the hyperplane at infinity is distinct from this hyperplane and contains no point of \(V'\) ; for example K can be taken to be infinite and \(K'\) finite.)
\end{enumerate}

¶ 18. Let \(\mathbf{E} = \mathbf{P}(\mathbf{V})\) be a projective plane over a field \(\mathbf{K}\) and \(u\) a bijective mapping of \(\mathbf{E}\) onto itself, arising when passing to the quotients from a bijective semi-linear mapping \(v\) of \(\mathbf{V}\) onto itself, relative to an automorphism \(\sigma\) of \(\mathbf{K}\) .

\begin{enumerate}
\def\labelenumi{(\alph{enumi})}
\item
  Show that the four following properties are equivalent: (α) for all \(x \in E\) , \(x, u(x)\) and \(u^{2}(x)\) lie on a straight line; (β) every line of E contains a point invariant under u; (γ) through every point of E there passes a line invariant under u; (δ) for every line D of E, the three line D, \(u(\mathbf{D})\) and \(u^{2}(\mathbf{D})\) have a common point. (Show first that (α) and (γ) are equivalent; deduce by duality (Exercise 12) that (β) and (δ) are equivalent; prove finally that (γ) implies (β) and deduce by duality that (β) implies (γ).)
\item
  Suppose that u has the properties stated in (a). Show that if there exists in E a line D invariant under u and containing only a single point invariant under u, u arises from a transvection v of V when passing to the quotients (§ 10, Exercise 11). (Show that every line invariant under u contains a; by considering a line not passing through a, show that there exists a line \(D_{0}\) passing through a and containing at least two points invariant under u; conclude that all the points of \(D_{0}\) are necessarily invariant under u.)
\item
  Suppose that u has the properties stated in (a). Show that if there exists in E a line E invariant under u and containing only two points invariant under u, u arises from a dilatation v of V when passing to the quotients ( \(\S\) 10, Exercise 11). (If a, b are the two points of D invariant under u, show that every line invariant under u passes through a or through b; then note that there exist at least two other points c, d distinct from a, b and invariant under u and therefore the line \(D_{cd}\) passes through a or b; conclude by proving that all the points of \(D_{cd}\) are invariant under u.)
\item
  Suppose that u has the properties stated in (a) and that every line of E invariant under u contains at least three distinct points invariant under u; then there exists in E a projective basis (Exercise 10) each point of which is invariant under u; conclude that there exists a basis \((e_{i})_{1\leq i\leq3}\) of V such that u arises from a semi-linear mapping \(v\) of \(\mathbf{V}\) into itself such that \(v(e_i) = e_i\) for \(1 \leqslant i \leqslant 3\) , when passing to the quotients. The set of points of \(\mathbf{E}\) invariant under \(u\) is then the projective plane \(\mathbf{P}(\mathbf{V}_0)\) , where \(\mathbf{V}_0\) is the vector space over the field \(\mathbf{K}_0\) of invariants of \(\sigma\) , generated by \(e_1, e_2, e_3\) .
\item
  Suppose henceforth that \(u\) satisfies the conditions of (a) and (d) and that neither \(u\) nor \(u^2\) is the identity. Show that there exists \(\gamma \in \mathbf{K}\) such that \(\gamma^\sigma = \gamma\) and
\end{enumerate}

\[
(\xi^ {\sigma} - \xi) ^ {\sigma} = \gamma (\xi^ {\sigma} - \xi)\tag{1}
\]

for all \(\xi \in \mathbf{K}\) (use conditions \((\alpha)\) of \((a)\) and the existence of \(\zeta \in \mathbf{K}\) such that \(\zeta^{\sigma} \neq \zeta\) ). Show that \(\gamma \neq -1\) and that, for all \(\xi \in \mathbf{K}\) such that \(\xi^{\sigma} \neq \xi\) ,

\[
(1 + \gamma) \xi^ {\sigma} \gamma = \gamma \xi (1 + \gamma)\tag{2}
\]

(apply (1) replacing \(\xi\) by \(\xi^2\) ); extend (2) to all \(\xi \in \mathbf{K}\) by noting that \(\xi = \eta - \zeta\) , where \(\eta^\sigma \neq \eta\) and \(\xi^\sigma \neq \zeta\) . Conclude that \(\gamma \neq 1\) , then deduce from (1) and (2) that

\[
\xi^ {\sigma} = (1 + \gamma) \xi (1 + \gamma) ^ {- 1}\tag{3}
\]

for all \(\xi \in K\) and that \(\gamma + \gamma^{-1}\) belongs to the centre of \(K\) . Obtain the converse. *Give an example where \(\gamma\) does not belong to the centre of \(K\) , but where \(\gamma + \gamma^{-1}\) is in this centre.*

*19. Let K be a commutative field, \(\tilde{K}\) the projective field obtained by adjoining to K a point at infinity (no. 9) and \(f\) and \(g\) two elements of K(X). Show that, writing \(h(\mathbf{X}) = f(g(\mathbf{X}))\) , if \(\tilde{f}, \tilde{g}, \tilde{h}\) are the canonical extensions of \(f, g, h\) to \(\tilde{\mathbf{K}}\) , then \(\tilde{h} = \tilde{f} \circ \tilde{g}_{\cdot *}\)

\BourbakiExercisesSection

\begin{enumerate}
\def\labelenumi{\arabic{enumi}.}
\tightlist
\item
  \begin{enumerate}
  \def\labelenumii{(\alph{enumii})}
  \tightlist
  \item
    Let E be a right A-module. On the additive group \(E^{r}\) (with \(r \geqslant 1\) ) an external law of composition is defined with \(\mathbf{M}_{r}(\mathbf{A})\) as set of operators, denoting by x.P, for every element \(x = (x_{i})_{1 \leqslant i \leqslant r}\) of \(E^{r}\) and every square matrix \(P = (\alpha_{ij}) \in \mathbf{M}_{r}(\mathbf{A})\) , the element \(y = (y_{i})\) of \(E^{r}\) such that
  \end{enumerate}
\end{enumerate}

\[
y _ {i} = \sum_ {j = 1} ^ {r} x _ {j} \alpha_ {j i} \quad (1 \leqslant i \leqslant r).
\]

This external law defines with the additive law on \(E^{r}\) a right \(\mathbf{M}_{r}(\mathbf{A})\) -module structure on \(E^{r}\) ; by restricting the ring of operators of the ring of scalar matrices \(I.\alpha\) , the product A-module structure on \(E^{r}\) is recovered. Show that for the A-module E to have a generating system of r elements, it is necessary and sufficient that the \(\mathbf{M}_{r}(\mathbf{A})\) -module \(E^{r}\) be monogenous.

\begin{enumerate}
\def\labelenumi{(\alph{enumi})}
\setcounter{enumi}{1}
\tightlist
\item
  Let \((x_{i})_{1\leqslant i\leqslant n}\) be a generating system of E and \((y_{j})_{1\leqslant j\leqslant m}\) a family of elements of E (resp. a generating system of E). Let \(z, z', z''\) be three elements of \(\mathbf{E}^{m + n}\) such that \(z_{i} = x_{i}\) for \(1\leqslant i\leqslant n\) , \(z_{n + j} = 0\) for \(1\leqslant j\leqslant m\) , \(z_{i}' = x_{i}\) for
\end{enumerate}

\(1 \leqslant i \leqslant n, z_{n+j}^{\prime} = y_{j}\) for \(1 \leqslant j \leqslant n, z_{i}^{\prime \prime} = 0\) for \(1 \leqslant i \leqslant n, z_{n+j}^{\prime \prime} = y_{j}\) for \(1 \leqslant j \leqslant n\) . Show that there exist two invertible matrices \(P, Q\) of \(\mathbf{M}_{n+m}(\mathbf{A})\) such that \(z' = z.P\) (resp. \(z'' = z.Q\) ).

\begin{enumerate}
\def\labelenumi{(\alph{enumi})}
\setcounter{enumi}{2}
\tightlist
\item
  If A is commutative, \(u_{P}: x \mapsto x \cdot P\) is an endomorphism of the A-module \(E^{r}\) and the mapping \(P \mapsto u_{P}\) is a homomorphism of the ring \(\mathbf{M}_{r}(\mathbf{A})\) into the ring \(\operatorname{End}_{\mathbf{A}}(\mathbf{E}^{r})\) . If E is a faithful A-module this homomorphism is injective.
\end{enumerate}

\begin{enumerate}
\def\labelenumi{\arabic{enumi}.}
\setcounter{enumi}{1}
\tightlist
\item
  \begin{enumerate}
  \def\labelenumii{(\alph{enumii})}
  \tightlist
  \item
    Let \(X\) be a square matrix over a ring \(\mathbf{A}\) , which can be written as an upper triangular matrix \((X_{ij})\) of matrices \((1 \leqslant i \leqslant p, 1 \leqslant j \leqslant p)\) . Show that if each of the square matrices \(X_{ii}\) \((1 \leqslant i \leqslant p)\) is invertible, so is \(X\) and \(X^{-1}\) can be written as an upper triangular matrix \((Y_{ij})\) \((1 \leqslant i \leqslant p, 1 \leqslant j \leqslant p)\) corresponding to the same partition of the indexing set. When \(\mathbf{A}\) is a field, prove that this sufficient condition for \(X\) to be invertible is also necessary.
  \end{enumerate}
\end{enumerate}

\begin{enumerate}
\def\labelenumi{(\alph{enumi})}
\setcounter{enumi}{1}
\tightlist
\item
  Give an example of a ring A and a matrix \(\begin{pmatrix} a & 1 \\ 0 & b \end{pmatrix}\) (with \(a \in \mathbf{A}, b \in \mathbf{A}\) ) which is invertible without either \(a\) or \(b\) being invertible in A and whose inverse \(\begin{pmatrix} a' & b' \\ c' & d' \end{pmatrix}\) is such that \(b' \neq 0\) and \(c' \neq 0\) (take A to be the endomorphism ring of an infinite-dimensional vector space).
\end{enumerate}

\begin{enumerate}
\def\labelenumi{\arabic{enumi}.}
\setcounter{enumi}{2}
\tightlist
\item
  Let A be the quotient ring Z/30Z. Show that in the matrix
\end{enumerate}

\[
\left( \begin{array}{c c c} 1 & 1 & - 1 \\ 0 & 2 & 3 \end{array} \right)
\]

over the ring A, the two rows are linearly independent but any two columns are linearly dependent.

\begin{enumerate}
\def\labelenumi{\arabic{enumi}.}
\setcounter{enumi}{3}
\tightlist
\item
  Let A be a ring, C its centre, B the ring of matrices \(\mathbf{M}_n(\mathbf{A}), \Delta\) the additive subgroup of A generated by the elements \(\alpha\beta - \beta\alpha\) for \(\alpha \in \mathbf{A}, \beta \in \mathbf{A}\) , and D the additive subgroup of B generated by the matrices \(XY - YX\) for \(X \in \mathbf{B}, Y \in \mathbf{B}\) ; \(\Delta\) and D are C-modules. For every matrix \(X = (\xi_{ij}) \in \mathbf{B}\) , let \(\theta(X)\) be the element \(\sum_{i=1}^{n} \bar{\xi}_{ij}\) of A/ \(\Delta\) , where, for all \(\alpha \in \mathbf{A}\) , \(\bar{\alpha}\) denotes the class of \(\alpha\) mod. \(\Delta\) . Show that \(\theta\) is a surjective homomorphism of B onto A/ \(\Delta\) , whose kernel is equal to D, so that B/D is isomorphic to A/ \(\Delta\) as a C-module. (Observe that D contains the matrices \(\alpha E_{ij}\) and the matrices \(\alpha(E_{ii} - E_{jj})\) for \(\alpha \in \mathbf{A}\) and \(i \neq j\) .)
\end{enumerate}

¶ 5. Let A be a ring, L, M, N any three indexing sets \(U = (a_{\lambda \mu})\) a matrix of \(\mathbf{A}^{\mathbf{L} \times \mathbf{M}}\) and \(V = (b_{\mu \nu})\) a matrix of \(\mathbf{A}^{\mathbf{M} \times \mathbf{N}}\) ; if, for every ordered pair \((\lambda, \nu) \in \mathbf{L} \times \mathbf{N}\) , the family \((a_{\lambda \mu} b_{\mu \nu})\) ( \(\mu \in \mathbf{M}\) ) has finite support, the element \(c_{\lambda \nu} = \sum_{\mu} a_{\lambda \mu} b_{\mu \nu}\) is defined and it is also said that the matrix \((c_{\lambda \nu})\) is the product \(UV\) of \(U\) by \(V\) . When the products \(UV'\) and \(UV''\) are defined, so is \(U(V' + V'')\)

and \(UV' + UV'' = U(V' + V'')\) ; is the converse true? Give an example of three infinite matrices, U, V, W such that the products UV, VW, \(U(VW)\) and \((UV)W\) are defined but \(U(VW) \neq (UV)W\) (take U to be a matrix with one row all of whose elements are equal to 1, W its transpose and V a matrix all of whose elements are 0, 1 or -1 and which has only a finite number of elements \(\neq 0\) in each row and each column. Determine by induction the elements of V such that UV = 0 but VW has a single element \(\neq 0\) .)

\begin{enumerate}
\def\labelenumi{\arabic{enumi}.}
\setcounter{enumi}{5}
\item
  Let X be a matrix with m rows and n columns over a field K; show that the rank \(\operatorname{rg}(X)\) is equal to the greatest of the ranks of the submatrices of X with an equal number of rows and columns. (Let \(\operatorname{rg}(X) = r\) ; if \(a_{1}, \ldots, a_{r}\) are r columns of X forming a free system in \(K_{d}^{m}\) and a basis of \(K_{d}^{m}\) is formed with these r vectors and m - r vectors of the canonical basis, show that the components of \(a_{1}, \ldots, a_{r}\) over the r other vectors of the canonical basis form a matrix of rank r.)
\item
  Let X be a matrix with m rows and n columns over a field K; if r is the rank of X, show that the rank of a submatrix with m rows and s columns, obtained by suppressing n - s columns of X, is \(\geqslant r + s - n\) .
\item
  Let \(X = (\alpha_{ij})\) be a matrix with \(m\) rows and \(n\) columns over a field \(K\) . For \(X\) to be of rank 1, it is necessary and sufficient that there exist in \(K\) a family \((\lambda_i)_{1 \leqslant i \leqslant m}\) of \(m\) elements not all zero and a family \((\mu_j)_{1 \leqslant j \leqslant n}\) of \(n\) elements not all zero such that \(\alpha_{ij} = \lambda_i \mu_j\) for every ordered pair of indices.
\item
  Let X, Y be two matrices with m rows and n columns over a field K; if there exist two square matrices \(P, P_{1}\) of order m and two square matrices \(Q, Q_{1}\) of order n such that Y = PXQ and \(X = P_{1}XQ_{1}\) , show that X and Y are equivalent.
\item
  Let \(X, X', Y, Y'\) be four square matrices of order \(n\) over a ring \(A\) , such that \(X\) is invertible. For there to exist two invertible square matrices \(P, Q\) of order \(n\) such that \(X' = PXQ\) and \(Y' = PYQ\) , it is necessary and sufficient that \(X'\) be invertible and that the matrices \(YX^{-1}\) and \(Y'X'^{-1}\) be similar.
\item
  Let E be a right vector space over a field K, of dimension \(\geqslant1\) , and H a hyperplane of E. Every endomorphism u of E leaving invariant each of the elements of H gives, on passing to the quotients, an endomorphism of the 1-dimensional quotient space E/H, an endomorphism which is therefore of the form \(\dot{x}\mapsto\dot{x}\mu(\dot{x})\) , where \(\mu(\dot{x})\in\mathrm{K}\) is such that \(\mu(\dot{x}\lambda)=\lambda^{-1}\mu(\dot{x})\lambda\) for \(\lambda\in K^{*}\) . An automorphism of E leaving invariant the elements of H is called a transvection of hyperplane H if the corresponding automorphism of E/H is the identity and a dilatation of hyperplane H otherwise; when u is a dilatation, the set of elements \(\mu(\dot{x})\) for \(\dot{x}\in E/H\) which is a class of conjugate elements in the multiplicative group \(K^{*}\) , is called the class of the dilatation u.
\end{enumerate}

\begin{enumerate}
\def\labelenumi{(\alph{enumi})}
\item
  Show that for every dilatation there exists one and only one supplementary line of H, invariant under the dilatation.
\item
  Let \(\phi\) be a linear form on \(\mathbf{E}\) such that \(\mathbf{H} = \overline{\phi}(0)\) ; show that for every transvection \(u\) of hyperplane \(\mathbf{H}\) there exists a unique vector \(a \in \mathbf{H}\) such that \(u(x) = x + a\phi(x)\) for all \(x \in \mathbf{E}\) . Let \(\Gamma(\mathbf{E}, \mathbf{H})\) be the subgroup of \(\mathbf{GL}(\mathbf{E})\) consisting of the automorphisms leaving invariant each element of \(\mathbf{H}\) ; show that the transvections of hyperplane \(\mathbf{H}\) form a normal commutative subgroup \(\Theta(\mathbf{E}, \mathbf{H})\) of \(\Gamma(\mathbf{E}, \mathbf{H})\) isomorphic to the additive group \(\mathbf{H}\) ; the quotient group \(\Gamma(\mathbf{E}, \mathbf{H}) / \Theta(\mathbf{E}, \mathbf{H})\) is isomorphic to the multiplicative group \(\mathbf{K}^*\) . For \(\Gamma(\mathbf{E}, \mathbf{H})\) to be commutative, it is necessary and sufficient that one of the two following conditions be satisfied: ( \(\alpha\) ) \(\mathbf{K}\) is a field with two elements; ( \(\beta\) ) \(\mathbf{K}\) is commutative and \(\dim \mathbf{E} = 1\) .
\item
  Suppose that E is finite-dimensional; show that for every transvection u there exists a basis of E such that the matrix of u with respect to this basis has all its diagonal elements equal to 1 and at most one other element \(\neq0\) .
\item
  Show that the centralizer (I, § 5, no. 3) of the group \(\Theta(E, H)\) in the group \(\mathbf{GL}(E)\) is the composition \(Z(E)\,\Theta(E, H) = \Theta(E, H)Z(E)\) of \(\Theta(E, H)\) and the centre \(Z(E)\) of \(\mathbf{GL}(E)\) (cf.~§ 1, Exercise 26). The only automorphisms belonging to this centralizer and leaving invariant at least one element \(\neq0\) of \(E\) are the transvections of the group \(\Theta(E, H)\) . If \(K\) has at least 3 elements, the centralizer of \(\Gamma(E, H)\) in \(\mathbf{GL}(E)\) is equal to \(Z(E)\) .
\item
  Show that the normalizers (I, § 5, no. 3) of \(\Theta(\mathbf{E},\mathbf{H})\) and \(\Gamma(\mathbf{E},\mathbf{H})\) in \(\mathbf{GL}(\mathbf{E})\) are both equal to the subgroup consisting of the automorphisms leaving H invariant.
\end{enumerate}

¶ 12. Let E be a right vector space of dimension ≥1 over a field K. Let F(E) denote the normal subgroup of GL(E) consisting of the automorphisms u such that the kernel of 1\_E - u is of finite codimension (hence F(E) = GL(E) if E is finite-dimensional). Let SL(E) denote the normal subgroup of GL(E) generated by all the transvections (Exercise 11); it is contained in F(E).

\begin{enumerate}
\def\labelenumi{(\alph{enumi})}
\item
  If \(\dim E \geqslant 2\) , show that for every ordered pair of non-zero vectors x, y of E, there exists a transvection or product of two transvections, which maps x to y (in other words \(\mathbf{SL}(\mathbf{E})\) operates transitively on \(E = \{0\}\) ).
\item
  Let V, W be two hyperplanes of E, \(\dot{x}_{0} = x_{0} + V\) a class mod. V distinct from V and \(\dot{y}_{0} = y_{0} + W\) a class mod. W distinct from W. Show that if \(\dim E \geqslant 2\) , there exists a transvection or product of two transvections, which maps V to W and \(\dot{x}_{0}\) to \(\dot{y}_{0}\) (consider first the case where V and W are distinct).
\item
  If \(\dim \mathbf{E} \geqslant 2\) , show that any two transvections distinct from the identity are conjugate in the group \(\mathbf{F}(\mathbf{E})\) .
\item
  If \(\dim \mathbf{E} \geqslant 3\) , show that any two transvections distinct from the identity are conjugate in the group \(\mathbf{SL}(\mathbf{E})\) (reduce it with the aid of (b) to the case where the hyperplanes of the two transvections are identical, then use (a)).
\item
  Suppose that \(\dim E = 2\) . For any two transvections to be conjugate in \(\mathbf{SL}(\mathbf{E})\) , it is necessary and sufficient that the subgroup \(\mathbf{Q}\) of \(\mathbf{K}^*\) generated by the squares of elements of \(\mathbf{K}^*\) be identical with \(\mathbf{K}^*\) . (If \(u\) is a transvection distinct from the identity, \(a\) a vector of \(\mathbf{E}\) which is not invariant under \(u\) and \(b = u(a) - a\) , show that, for every transvection \(u'\) conjugate to \(u\) in \(\mathbf{SL}(\mathbf{E})\) , \(u'(a) - a = a\lambda + b\mu\) , with \(\mu = 0\) or \(\mu \in \mathbf{Q}\) ; for this use the fact that in a group \(\mathbf{G}\) every product \(sts^{-1}t\) is a product of squares.)
\item
  Suppose that \(\dim E \geqslant 2\) . For two dilatations \(u, u'\) to be such that there exists a \(v \in \mathbf{SL}(\mathbf{E})\) such that \(u' = vuv^{-1}\) , it is necessary and sufficient that the classes (Exercise 11) of \(u\) and \(u'\) be the same (use \((b)\) ). Deduce that if the class of a dilatation is contained in the commutator group of \(\mathbf{K}^*\) , this dilatation belongs to the group \(\mathbf{SL}(\mathbf{E})\) (observe that \((vuv^{-1})u^{-1} = v(uv^{-1}u^{-1}))\) .
\end{enumerate}

¶ 13. Let E be a right vector space of dimension ≥2 and H₀ a hyperplane of E.

\begin{enumerate}
\def\labelenumi{(\alph{enumi})}
\item
  Show that every automorphism \(u\) belonging to \(\mathbf{F}(\mathbf{E})\) (Exercise 12) is the product of an automorphism of \(\mathbf{SL}(\mathbf{E})\) and possibly a dilatation of the hyperplane \(\mathbf{H}_0\) (proceed by induction on the codimension of the kernel of \(1_{\mathbf{E}} - u\) , using Exercises 12(a) and 12(f)).
\item
  Show that \(\mathbf{SL}(\mathbf{E})\) contains the commutator group of \(\mathrm{F(E)}\) and is identical with this group except when \(\mathbf{E}\) is a space of dimension 2 over the field with 2 elements. (To show that \(\mathbf{SL}(\mathbf{E})\) contains the commutator group of \(\mathrm{F(E)}\) , use (a) and Exercise \(12(f)\) . To see that \(\mathbf{SL}(\mathbf{E})\) is contained in the commutator group of \(\mathrm{F(E)}\) , except in the case indicated, show that for every homomorphism of \(\mathrm{F(E)}\) into a commutative group, the image of a transvection is the identity element, using Exercises \(11(b)\) and \(12(c)\) .)
\item
  Show that \(\mathbf{SL}(\mathbf{E})\) is equal to its commutator group except when \(\dim \mathbf{E} = 2\) and \(\mathbf{K}\) has 2 or 3 elements. (If \(\dim \mathbf{E} \geqslant 3\) , note that every transvection \(u\) can be written as \(vwv^{-1}w^{-1}\) , where \(v\) is a transvection and \(w \in \mathbf{SL}(\mathbf{E})\) , using Exercise 12(d). If \(\dim \mathbf{E} = 2\) , note first that every automorphism \(u\) whose matrix with respect to a basis of \(\mathbf{E}\) is of the form \(\begin{pmatrix} \lambda & 0 \\ 0 & \lambda^{-1} \end{pmatrix}\) is a product of transvections, arguing as in (a); then consider the commutator \(uvu^{-1}v^{-1}\) , where \(v\) is the transvection whose matrix with respect to the same basis is \(\begin{pmatrix} 1 & 1 \\ 0 & 1 \end{pmatrix}\) .)
\end{enumerate}

¶ 14. Let E be a vector space of dimension ≥2 over a field K.

\begin{enumerate}
\def\labelenumi{(\alph{enumi})}
\item
  Show that the projective group PGL(E) is canonically isomorphic to the quotient of the linear group GL(E) by its centre (isomorphic to the multiplicative group of the centre of K, cf.~§ 1, Exercise 26).
\item
  Let \(\mathbf{PSL}(\mathbf{E})\) denote the canonical image of the group \(\mathbf{SL}(\mathbf{E})\) in \(\mathbf{PGL}(\mathbf{E})\) ; it is a normal subgroup of \(\mathbf{PGL}(\mathbf{E})\) , containing the commutator group of \(\mathbf{PGL}(\mathbf{E})\) when \(\mathbf{E}\) is finite-dimensional. Show that \(\mathbf{PSL}(\mathbf{E})\) is a doubly transitive (cf.~I, § 5, Exercise 14) group of permutations of \(\mathbf{P}(\mathbf{E})\) .
\item
  Let \(a \neq 0\) be an element of \(\mathbf{E}\) and let \(e = \pi(a) \in \mathbf{P}(\mathbf{E})\) (in the notation of §9, no. 6). Show that the transvections of the form \(x \mapsto x + a\phi(x)\) (where \(\phi \in \mathbf{E}^*\) is such that \(\phi(a) = 0\) ) constitute a commutative subgroup of \(\mathbf{SL}(\mathbf{E})\) ; if \(\Gamma_e\) is the image of this subgroup in \(\mathbf{PSL}(\mathbf{E})\) , \(\Gamma_e\) is a normal subgroup of the subgroup \(\Phi_e\) of \(\mathbf{PSL}(\mathbf{E})\) leaving \(e\) invariant. Show also that the union of the subgroups conjugate to \(\Gamma_e\) in \(\mathbf{PSL}(\mathbf{E})\) generates \(\mathbf{PSL}(\mathbf{E})\) .
\item
  Let \(\Sigma\) be a primitive group (I, § 5, Exercise 13) of permutations of a set F which is equal to its commutator group. Suppose that there exists an element \(c \in \mathbf{F}\) such that the subgroup \(\Phi_c\) of \(\Sigma\) leaving \(c\) invariant contains a commutative normal subgroup \(\Gamma_c\) such that the union of the subgroups conjugate to \(\Gamma_c\) in \(\Sigma\) generates \(\Sigma\) . Show that under these conditions \(\Sigma\) is simple. (Let \(\Delta\) be a normal subgroup of \(\Sigma\) distinct from the identity; using I, § 5, Exercise 17, show that \(\Sigma = \Delta\) . \(\Phi_c = \Delta\) . \(\Gamma_c\) and, using the fact that \(\Gamma_c\) is commutative, deduce that every commutator in \(\Sigma\) is contained in \(\Delta\) .)
\item
  Deduce from (c), (d) and Exercise 13(c) that the group \(\mathbf{PSL}(\mathbf{E})\) is simple except when \(\dim \mathbf{E} = 2\) and \(\mathbf{K}\) has 2 or 3 elements.
\end{enumerate}

\(^{*}(f)\) If \(\dim E = n + 1\) and K is a field with q elements, show that the order of PGL(E) is

\[
(q ^ {n + 1} - 1) (q ^ {n + 1} - q) \dots (q ^ {n + 1} - q ^ {n - 1}) q ^ {n}
\]

(use (a), Exercise 9(a) of § 9, and the fact that K is necessarily commutative (V, § 11, Exercise 14 and VIII, § 11, no. 1, Theorem 1)).*

\begin{enumerate}
\def\labelenumi{(\alph{enumi})}
\setcounter{enumi}{6}
\item
  Show that if \(\dim \mathbf{E} = 2\) and \(\mathbf{K}\) is a field with 2 (resp. 3) elements, \(\mathbf{PSL}(\mathbf{E})\) is isomorphic to the symmetric group \(\mathfrak{S}_3\) (resp. the alternating group \(\mathfrak{A}_4\) ). (Considering \(\mathbf{PSL}(\mathbf{E})\) as a group of permutations of \(\mathbf{P}(\mathbf{E})\) , note that it contains every transposition (resp. every cycle ( \(a b c\) ) and use \(f\) ).)
\item
  Show that, except in the cases considered in (g), every normal subgroup of \(\mathbf{SL}(\mathbf{E})\) distinct from \(\mathbf{SL}(\mathbf{E})\) is contained in the centre of \(\mathbf{SL}(\mathbf{E})\) (use (d) and the fact that every transvection is a commutator in \(\mathbf{SL}(\mathbf{E})\) (Exercise 13(c))).
\end{enumerate}

¶ 15. Let A be a ring; for every integer \(n \geqslant 1\) , \(\mathbf{A}^{n}\) is canonically identified with the submodule of \(\mathbf{A}^{(\mathbf{N})}\) consisting of the elements whose coordinates of index \(\geqslant n\) are 0 and \(\mathbf{A}_{n}^{\prime}\) denotes the complementary submodule consisting of the elements whose coordinates of index \(<n\) are 0. Every endomorphism \(u \in \mathbf{GL}_{n}(\mathbf{A})\) is identified with the automorphism of \(\mathbf{A}^{(\mathbf{N})}\) whose restriction to \(\mathbf{A}^{n}\) is \(u\) and whose restriction to \(\mathbf{A}_{n}^{\prime}\) is the identity, so that \(\mathbf{GL}_{n}(\mathbf{A})\) is identified with a subgroup of \(\mathbf{GL}(\mathbf{A}^{(\mathbf{N})})\) ; let F denote the subgroup of \(\mathbf{GL}(\mathbf{A}^{(\mathbf{N})})\) the union of the \(\mathbf{GL}_{n}(\mathbf{A})\) , in other words the subgroup leaving invariant the elements of the canonical basis of \(\mathbf{A}^{(\mathbf{N})}\) except for a finite number of them (cf.~Exercise 12).

Let \(\mathbf{T}_n\) be the subgroup of \(\mathbf{GL}_n(\mathbf{A})\) generated by the \(\mathrm{B}_{ij}(\lambda)\) for \(0 \leqslant i < n\) , \(0 \leqslant j < n\) , \(i \neq j\) , \(\lambda \in \mathbf{A}\) (no. 13) and let \(\mathbf{T}\) be the subgroup of \(\mathbf{F}\) the union of the \(\mathbf{T}_n\) .

\begin{enumerate}
\def\labelenumi{(\alph{enumi})}
\item
  Show that for \(n \geqslant 3\) , \(T_n\) is equal to its commutator group.
\item
  Let \(a, b\) be two invertible elements of \(A\) . Show that
\end{enumerate}

\[
\left( \begin{array}{c c} a b & 0 \\ 0 & 1 \end{array} \right) = P \left( \begin{array}{c c} a & 0 \\ 0 & b \end{array} \right) Q \quad \text { and } \quad \left( \begin{array}{c c} a & 0 \\ 0 & b \end{array} \right) = P ^ {\prime} \left( \begin{array}{c c} b & 0 \\ 0 & a \end{array} \right) Q ^ {\prime}
\]

where \(P, Q, P', Q'\) belong to \(T_2\) . Deduce that the matrix \(\begin{pmatrix} aba^{-1}b^{-1} & 0 \\ 0 & 1\end{pmatrix}\) belongs to \(T_2\) .

\begin{enumerate}
\def\labelenumi{(\alph{enumi})}
\setcounter{enumi}{2}
\tightlist
\item
  Deduce from (b) that the commutator group of \(\mathbf{GL}_{n}(\mathbf{A})\) is contained in \(T_{2n}\) (replace A by \(\mathbf{M}_{n}(\mathbf{A})\) in (b)). Conclude (using (a)) that the commutator group of F is equal to T.
\end{enumerate}

\begin{enumerate}
\def\labelenumi{\arabic{enumi}.}
\setcounter{enumi}{15}
\tightlist
\item
  \begin{enumerate}
  \def\labelenumii{(\alph{enumii})}
  \tightlist
  \item
    Let \(U\) be a square matrix of order \(n\) over the field \(\mathbf{Q}\) all of whose non-diagonal elements are equal to the same rational number \(r > 0\) and whose diagonal \((d_1, \ldots, d_n)\) is such that \(d_i \geqslant r\) for all \(i\) and \(d_i > r\) except perhaps for one value of \(i\) . Show that \(U\) is invertible (if \(x = (x_i)_{1 \leqslant i \leqslant n}\) is a solution of the equation \(U.x = 0\) , show that the \(x_i\) all have the same sign and deduce that they are all zero).
  \end{enumerate}
\end{enumerate}

\begin{enumerate}
\def\labelenumi{(\alph{enumi})}
\setcounter{enumi}{1}
\tightlist
\item
  Let E be a finite set with m elements denoted by \(a_{i}\) ( \(1 \leqslant i \leqslant m\) ) and \((\mathbf{A}_{j})_{1 \leqslant j \leqslant n}\) a family of n subsets of E, distinct from one another; suppose that \(\text{Card}(\mathbf{A}_{i} \cap \mathbf{A}_{j}) = r \geqslant 1\) for every ordered pair \((i, j)\) of distinct elements of \([1, n]\) . Show that necessarily \(n \leqslant m\) . (Let \(V = (c_{ij})\) be the matrix of type \((m, n)\) such that \(c_{ij} = 1\) when \(a_{i} \in A_{j}\) , \(c_{ij} = 0\) otherwise. Apply (a) to the matrix \(^{t}V\) . V of order n.)
\end{enumerate}

¶ 17. Let K be a field (commutative or otherwise), E an n-dimensional right vector space over K and u an endomorphism of E.

\begin{enumerate}
\def\labelenumi{(\alph{enumi})}
\item
  Show that if \(n > 1\) and \(u\) is not a central homothety, there exists a basis of \(E\) such that the matrix of \(u\) with respect to this basis is of the form \((\alpha_{ij})\) with \(\alpha_{ij} = 0\) for \(j > i + 1\) and \(\sum_{i=1}^{n-1} \alpha_{i,i+1} = 1\) (argue by induction on \(n\) ).
\item
  Show that if \(n > 2\) , there exists a basis of \(\mathbf{E}\) such that the matrix of \(u\) with respect to this basis is of the form \((\alpha_{ij})\) with \(\alpha_{ij} = 0\) for \(j > i + 1\) and
\end{enumerate}

\(\sum_{i=1}^{n-1} \alpha_{i, i+1} = 0\) . (Argue by induction on \(n\) using \((a)\) , noting that it amounts to the same to prove the proposition for \(u\) or for \(u + \gamma 1_{E}\) , where \(\gamma\) belongs to the centre of \(K\) , and studying separately the case where \(u\) is of rank 1 or 2.)

\begin{enumerate}
\def\labelenumi{(\alph{enumi})}
\setcounter{enumi}{2}
\item
  Let \(A = (\alpha_{ij})\) be a matrix with the properties described in (b). Also let \(S\) be the matrix of order \(n(\varepsilon_{ij})\) such that \(\varepsilon_{ij} = 0\) except when \(i = j + 1\) , in which case \(\varepsilon_{j+1,j} = 1\) . Show that there exists a matrix \(B = (\beta_{ij})\) of order \(n\) such that \(A = \lambda E_{nn} + (BS - SB)\) with \(\lambda \in \mathbf{K}\) . (Take \(B\) such that \(\beta_{i1} = 0\) for \(1 \leqslant i \leqslant n\) and \(\beta_{ij} = 0\) for \(j > i + 2\) .)
\item
  Suppose that K is commutative. Deduce from (c) that every endomorphism u of E such that \(\operatorname{Tr}(u) = 0\) can be written as vw - wv, where v and w are endomorphisms of E.
\end{enumerate}

\BourbakiExercisesSection

\begin{enumerate}
\def\labelenumi{\arabic{enumi}.}
\tightlist
\item
  Let \(\Delta\) be a commutative monoid, written additively, with its identity element denoted by 0. Let \(\mathbf{Z}^{(\Delta)}\) be the algebra of \(\Delta\) over \(\mathbf{Z}\) (III, §2, no. 6), that is a commutative ring, which is a free \(\mathbf{Z}\) -module with a basis \((\mathrm{X}^{\sigma})_{\sigma \in \Delta}\) and the multiplication \((a\mathrm{X}^{\sigma})(b\mathrm{X}^{\tau}) = ab\mathrm{X}^{\sigma +\tau}\) . For every \(\mathbf{Z}\) -module, we write \(\mathrm{E}^{(\Delta)} = \mathrm{E} \otimes_{\mathbf{Z}} \mathbf{Z}^{(\Delta)}\) , every element of \(\mathrm{E}^{(\Delta)}\) hence being written uniquely in the form \(\sum_{\sigma \in \Delta} z_{\sigma} \otimes \mathrm{X}^{\sigma}\) , with \(z_{\sigma} \in \mathrm{E}\) , the family \((z_{\sigma})\) having finite support. For every \(\mathbf{Z}\) -module homomorphism \(f: \mathrm{E} \to \mathrm{E}'\) , let \(f^{(\Delta)}\) denote the \(\mathbf{Z}^{(\Delta)}\) -module homomorphism \(f \otimes 1: \mathrm{E}^{(\Delta)} \to \mathrm{E}'^{(\Delta)}\) . Similarly, if \(\Delta'\) is another additive monoid with its identity element denoted by 0 and if \(\alpha: \Delta \to \Delta'\) is a monoid homomorphism, \(\alpha(\mathrm{E})\) denotes the homomorphism \(\mathrm{E}^{(\Delta)} \to \mathrm{E}^{(\Delta')}\) such that \(\alpha(\mathrm{E})(z_{\sigma} \otimes \mathrm{X}^{\sigma}) = z_{\sigma} \otimes \mathrm{X}^{\alpha(\sigma)}\) . \((\mathrm{E}^{(\Delta)})^{(\Delta')}\) is identified with \(\mathrm{E}^{(\Delta \times \Delta')}\) by associating with \(\sum_{\sigma' \in \Delta'} \left( \sum_{\sigma \in \Delta} z_{\sigma, \sigma'} \otimes \mathrm{X}^{\sigma} \right) \otimes \mathrm{X}^{\sigma'}\) the element \(\sum_{\sigma, \sigma'} z_{\sigma, \sigma'} \otimes \mathrm{X}^{(\sigma, \sigma')}.\) Finally, let \(\varepsilon(\mathrm{E})\) denote the homomorphism \(\mathrm{E}^{(\Delta)} \to \mathrm{E}\) mapping \(\sum_{\sigma} z_{\sigma} \otimes \mathrm{X}^{\sigma}\) to \(\sum_{\sigma} z_{\sigma}\)
\end{enumerate}

\begin{enumerate}
\def\labelenumi{(\alph{enumi})}
\tightlist
\item
  Let E be a Z-module, \((\mathrm{E}_{\sigma})_{\sigma\in\Delta}\) a graduation on E of type \(\Delta\) and, for all \(\sigma\in\Delta\) , let \(p_{\sigma}\) be the projector \(E\to E_{\sigma}\) corresponding to the decomposition of E as the direct sum of the \(E_{\sigma}\) . For all \(x\in E\) , let \(\phi_{\mathbf{E}}(x)=\sum_{\sigma\in\Delta}p_{\sigma}(x)\otimes X^{\sigma}\) . Show that the homomorphism \(\phi_{E}\colon E\to E^{(\Delta)}\) thus defined has the following properties:
\end{enumerate}

\[
\varepsilon (\mathbf {E}) \circ \phi_ {\mathbf {E}} = 1 _ {\mathbf {E}}\tag{1}
\]

\[
\delta (\mathbf {E}) \circ \phi_ {\mathbf {E}} = \phi_ {\mathbf {E}} ^ {(\Delta)} \circ \phi_ {\mathbf {E}},\tag{2}
\]

where \(\delta: \Delta \to \Delta \times \Delta\) is the diagonal mapping.

Show that a bijection of the set of graduations on E of type \(\Delta\) is thus defined on to the set of homomorphisms of E into \(\mathrm{E}^{(\Delta)}\) satisfying conditions (1) and (2).

\begin{enumerate}
\def\labelenumi{(\alph{enumi})}
\setcounter{enumi}{1}
\item
  Let E, F be two graded Z-modules of type \(\Delta\) . For a homomorphism \(f: E \to F\) to be graded of degree 0, it is necessary and sufficient that \(\phi_{F} \circ f = f^{(\Delta)} \circ \phi_{E}\) . For a submodule \(E'\) of E to be graded, it is necessary and sufficient that \(\phi_{E}(E') \subset E'^{(\Delta)}\) .
\item
  Let \(\mathbf{A}\) be a ring; \(\mathbf{A}^{(\Delta)} = \mathbf{A} \otimes_{\mathbf{Z}} \mathbf{Z}^{(\Delta)}\) is given a ring structure such that \((a \otimes \mathbf{X}^{\sigma})(b \otimes \mathbf{X}^{\tau}) = (ab) \otimes \mathbf{X}^{\sigma + \tau}\) . Let \((\mathbf{A}_{\sigma})_{\sigma \in \Delta}\) be a graduation of the additive
\end{enumerate}

group A; for this graduation to be compatible with the ring structure on A and \(1 \in A_{0}\) , it is necessary and sufficient that \(\phi_{A}: A \to A^{(\Delta)}\) be a ring homomorphism. State and prove an analogous result for graded A-modules.

\BourbakiAppendixExercises

\begin{enumerate}
\def\labelenumi{\arabic{enumi}.}
\tightlist
\item
  Let A be a pseudo-ring and M a left pseudomodule over A such that for all \(z \neq 0\) in M, \(z \in Az\) . Show that if x, y are two elements of M such that \(\text{Ann}(x) \subset \text{Ann}(y)\) in A, there exists an endomorphism u of M such that \(u(x) = y\) (prove that the relation bx = ax for a, b in A implies by = ay).
\end{enumerate}

\chapter{Tensor Algebras, Exterior Algebras, Symmetric Algebras}

Recall the exponential notation introduced in Chapter I, of which we shall make frequent use (I, § 7, no. 8):

Let \((x_{\lambda})_{\lambda \in L}\) be a family of pairwise permutable elements of a ring \(\mathbf{A}\) ; for every mapping \(\alpha: \mathbf{L} \to \mathbf{N}\) of finite support we shall write

\[
x ^ {\alpha} = \prod_ {\lambda \in L} x _ {\lambda} ^ {\alpha (\lambda)}.
\]

If \(\beta\) is another mapping of \(\mathbf{L}\) into \(\mathbf{N}\) of finite support, \(\alpha + \beta\) denotes the mapping

\[
\lambda \mapsto \alpha (\lambda) + \beta (\lambda)
\]

of L into N; with this law of composition the set \(\mathbf{N}^{(L)}\) of mappings of L into N of finite support is the free commutative monoid derived from L and

\[
x ^ {\alpha} x ^ {\beta} = x ^ {\alpha + \beta},
\]

For all \(\alpha \in \mathbf{N}^{(L)}\) , we write \(|\alpha| = \sum_{\lambda \in L} \alpha(\lambda) \in \mathbf{N}\) ; then \(|\alpha + \beta| = |\alpha| + |\beta|\) ; \(|\alpha|\) is called the order of the ``multiindex'' \(\alpha\) . For all \(\lambda \in L\) , let \(\delta_{\lambda}\) denote the element of \(\mathbf{N}^{(L)}\) such that \(\delta_{\lambda}(\lambda) = 1\) , \(\delta_{\lambda}(\mu) = 0\) for \(\mu \neq \lambda\) (Kronecker index); the \(\delta_{\lambda}\) for \(\lambda \in L\) are the only elements of \(\mathbf{N}^{(L)}\) of order 1. \(\mathbf{N}^{(L)}\) is given the ordering induced by the product ordering on \(\mathbf{N}^{\mathrm{L}}\) , so that the relation \(\alpha \leqslant \beta\) is equivalent to `` \(\alpha(\lambda) \leqslant \beta(\lambda)\) for all \(\lambda \in L\) ''; then the multiindex \(\lambda \mapsto \beta(\lambda) - \alpha(\lambda)\) is denoted by \(\beta - \alpha\) , so that it is the unique multiindex such that \(\alpha + (\beta - \alpha) = \beta\) . For all \(\alpha \in \mathbf{N}^{(L)}\) , there are only a finite number of multiindices \(\beta \leqslant \alpha\) ; the \(\delta_{\lambda}\) are the minimal elements of the set \(\mathbf{N}^{(L)} = \{0\}\) ; the relation \(\alpha \leqslant \beta\) implies \(|\alpha| \leqslant |\beta|\) and if both \(\alpha \leqslant \beta\) and \(|\alpha| = |\beta|\) , then \(\alpha = \beta\) . Finally, we write \(\alpha! = \prod_{\lambda \in L} (\alpha(\lambda))!\) , which is meaningful since \(0! = 1\) .

From § 4 to § 8 inclusive, A denotes a commutative ring and, unless otherwise stated, the algebras considered are assumed to be associative and unital and the algebra homomorphisms are assumed to be unital.

\section{ALGEBRAS}
